\documentclass[UTF8]{ctexbook}

% 支持从项目根目录或本文件所在目录编译。
\makeatletter
\def\input@path{{../}}
\makeatother
\usepackage{researchnotes}

\title{初等数学笔记}
\author{}
\date{}

\begin{document}
\maketitle
\tableofcontents

\chapter*{绪论：研究对象、课程总纲与学习路线}
\addcontentsline{toc}{chapter}{绪论：研究对象、课程总纲与学习路线}

\section*{初等数学研究什么}
\keyterm{初等数学}（elementary mathematics）以数、式、方程、函数和几何图形为基本对象，研究数量关系、空间关系以及有限过程中的规律。它既包含算术、初等代数和欧氏几何，也包含三角学、解析几何、初步组合计数与概率统计。本书将这些内容组织成一条连续的学习路线：从正确表达问题开始，用运算和推理建立关系，再根据关系计算、证明、估计或建模。

例如，“一定周长的长方形怎样使面积最大”可以用代数式、二次函数或不等式处理；“怎样测量无法直接到达的两点间距离”可以用相似三角形、三角函数或坐标处理。同一个问题的不同解法揭示不同结构，学习目标因而不仅是得到答案，还包括说明变量范围、变形依据、结论条件与方法之间的联系。

\section*{全书定位与当前版本}
本书面向希望系统补足初等数学基础、并准备学习大学数学的读者。先修要求是整数和分数的基本运算、常用度量单位及对简单几何图形的直观认识；不预设微积分、线性代数或抽象代数。基础内容覆盖中学数学的主要知识，并以本科生需要的证明意识、结构理解和建模表达加以统一；较深的凸性、图论与极限讨论可作为选读。

本书按本科阶段基础巩固与知识衔接的需要编写，正文覆盖全部二十四章，以定义、主要结论、关键证明和代表例题建立主线，章末配置习题与答案要点。论述尽量紧凑，但保留变量范围、等价变形依据和取等条件。实数系的基本结构、欧氏几何公理及部分几何基础定理作为先修事实使用；实指数延拓、连续函数介值性质和球面面积等未完整证明的结论在使用处明确说明边界。

\section*{课程总纲}
全书分六篇、二十四章，按照知识依赖安排。以下总纲概括课程结构，随后各篇展开正文。
\begin{description}
  \item[第一篇：数学语言、数与式]
  第\ref{chap:em-language}章，集合、逻辑与证明；第\ref{chap:em-numbers}章，数系、运算与近似；第\ref{chap:em-integers}章，整除、素数与同余初步；第\ref{chap:em-expressions}章，代数式与多项式。建立表达、计算和论证的共同基础。
  \item[第二篇：方程、不等式与函数]
  第\ref{chap:em-equations}章，方程与方程组；第\ref{chap:em-solving-inequalities}章，不等式与不等式组；第\ref{chap:em-functions}章，函数与图像；第\ref{chap:em-exponential}章，幂、指数与对数。以解集和定义域贯穿代数变形与函数研究。
  \item[第三篇：数列、估计与有限方法]
  第\ref{chap:em-sequences}章，数列、递推与求和；第\ref{chap:em-inequality-methods}章，初等不等式与最值；第\ref{chap:em-counting}章，计数原理与二项式定理；第\ref{chap:em-discrete}章，有限问题中的证明方法。联系规律发现、严格证明与有效计算。
  \item[第四篇：平面与空间几何]
  第\ref{chap:em-triangles}章，直线、三角形与相似；第\ref{chap:em-circles}章，四边形与圆；第\ref{chap:em-transformations}章，几何变换、轨迹与作图；第\ref{chap:em-solid}章，立体几何与度量。以欧氏几何中的位置关系、全等、相似和度量为主线。
  \item[第五篇：三角、坐标与复数]
  第\ref{chap:em-trigonometry}章，角与三角函数；第\ref{chap:em-triangle-solving}章，三角恒等变换、方程与解三角形；第\ref{chap:em-vectors}章，平面向量、直线与圆的坐标方法；第\ref{chap:em-conics}章，圆锥曲线；第\ref{chap:em-complex}章，复数及其几何意义。用代数语言统一描述长度、方向、位置与旋转。
  \item[第六篇：随机、数据与综合应用]
  第\ref{chap:em-probability}章，有限概率模型；第\ref{chap:em-statistics}章，统计描述与数据解释；第\ref{chap:em-modeling}章，数学建模、综合方法与后续学习。把确定性工具用于带有随机性或现实约束的问题。
\end{description}
附录\ref{app:em-notation}整理符号、公式与方法索引，附录\ref{app:em-study}提供分层自测及专题阅读路径。附录用于定位正文中的定义、条件与推导，不能替代正文学习。

\section*{学习路径与章节依赖}
基础学习沿前两篇顺序推进，随后可分两条线阅读：代数线为数列、不等式、计数，再到概率与统计；几何线为三角形、圆、几何变换、立体几何，再到三角与坐标。第五篇使用第二篇的函数语言及第四篇的几何基础，复数的三角形式还依赖三角恒等式。第\ref{chap:em-discrete}章的提高题、第\ref{chap:em-inequality-methods}章的凸性、第\ref{chap:em-conics}章的综合轨迹题可以在第一遍学习时暂缓。

已有专题笔记中的均值、伯努利不等式与 Jensen 不等式，分别归入第\ref{chap:em-inequality-methods}章的主干或选读内容；裂项求和与双阶乘分别归入第\ref{chap:em-sequences}章和第\ref{chap:em-counting}章；三角公式与解三角形对应第\ref{chap:em-trigonometry}、\ref{chap:em-triangle-solving}章；圆和球的度量对应第\ref{chap:em-circles}、\ref{chap:em-solid}章。正文按本书的先修顺序组织必要内容，较深的专题论证作为补充阅读，不要求读者预先读过这些独立笔记。

\section*{符号、论证与练习约定}
记自然数集为 $\mathbb N=\{0,1,2,\ldots\}$，正整数集为 $\mathbb N_{>0}$，整数、有理数、实数与复数集分别为 $\mathbb Z,\mathbb Q,\mathbb R,\mathbb C$；复数在第\ref{chap:em-complex}章正式引入。未知数通常记为 $x,y$，数列项记为 $a_n$，函数记为 $f$。同一字母在不同章节可能承担不同角色，使用时重新说明；几何中的 $A,B,C$ 通常表示点，三角形内角则写为 $\angle A,\angle B,\angle C$，需要简记时作明确约定。

等号表示相等，$\Longleftrightarrow$ 表示等价，$\Longrightarrow$ 表示单向推出，$\approx$ 表示带有上下文精度要求的近似。方程和不等式的变形始终保留定义域；除以可能为零的量必须分类，平方或开方必须检查符号。几何图形用于理解关系，不能代替证明；随机试验的样本空间和等可能性必须说明；实际数据与教学自拟数据严格区分。

章末习题结合各章内容安排概念辨析、基本计算、证明与反例或综合应用，附录提供分阶段检查任务。基础题训练正确使用定义，提高题训练选择方法，综合题训练在不同表示之间转换。代表例题给出关键依据与验算，章末习题附答案或提示；数值计算和作图用于探索与核验，不能据有限个例子断言一般规律。

\part{数学语言、数与式}

\chapter{集合、逻辑与证明}
\label{chap:em-language}
数学中的困难往往不只在计算，还在于条件是否清楚、推理是否有效。集合描述研究范围，逻辑说明条件之间的关系，证明则把已知事实连接到待证结论。本章采用朴素集合语言，不讨论集合论的公理化基础。

\section{集合、元素与常用表示}
\begin{definition}
\keyterm{集合}（set）是由确定对象组成的整体，对象称为元素。$x\in A$ 表示 $x$ 是集合 $A$ 的元素；$A\subseteq B$ 表示 $A$ 的每个元素都属于 $B$。若 $A\subseteq B$ 且 $B\subseteq A$，则称两集合相等。没有元素的集合记为 $\varnothing$。
\end{definition}
集合不记录排列次序和重复次数，故 $\{1,2,2\}=\{2,1\}$。描述法写成 $A=\{x\in D:P(x)\}$，其中 $D$ 是对象范围，$P(x)$ 是筛选条件。证明两个集合相等，通常分别证明两方向的包含。空集是每个集合的子集：若这一结论失败，就必须在空集中找到一个不属于目标集合的元素，而空集没有元素。

实数区间 $(a,b)=\{x:a<x<b\}$ 不含端点，$[a,b]=\{x:a\le x\le b\}$ 包含端点；半开区间相应处理。符号 $\pm\infty$ 不表示实数端点，因此区间在无穷一侧总用圆括号。必须区分元素关系 $1\in\{1,2\}$ 与包含关系 $\{1\}\subseteq\{1,2\}$。

\section{集合运算与分类表达}
在固定全集 $U$ 中，定义并集 $A\cup B=\{x:x\in A\text{ 或 }x\in B\}$、交集 $A\cap B=\{x:x\in A\text{ 且 }x\in B\}$、差集 $A\setminus B=\{x:x\in A,\ x\notin B\}$ 和补集 $A^c=U\setminus A$。这里“或”允许两者同时成立；“恰有一个成立”对应 $(A\setminus B)\cup(B\setminus A)$。

\begin{proposition}[德摩根律]
对 $U$ 的任意子集 $A,B$，有 $(A\cup B)^c=A^c\cap B^c$ 及 $(A\cap B)^c=A^c\cup B^c$。
\end{proposition}
\begin{proof}
元素 $x$ 不属于 $A\cup B$，当且仅当它既不属于 $A$ 也不属于 $B$，这正是属于 $A^c\cap B^c$ 的条件。第二式中“不同时属于两者”等价于“至少不属于其中之一”，同理得到结论。
\end{proof}
把一个集合写成若干子集的并称为分类；若各部分两两不交且覆盖原集合，称为分割。分类讨论必须覆盖全部情形；计数时若各类有交叠，还需修正重复。

\section{命题、量词与条件关系}
\keyterm{命题}（proposition）是具有确定真假值的陈述。$P\Rightarrow Q$ 表示只要 $P$ 成立，$Q$ 就成立，称 $P$ 是 $Q$ 的充分条件、$Q$ 是 $P$ 的必要条件；双向均成立写为 $P\Leftrightarrow Q$。例如对实数 $x$，$x>0$ 能推出 $x^2>0$，逆向不成立；正确的等价关系是 $x^2>0\Leftrightarrow x\ne0$。逆否命题 $\neg Q\Rightarrow\neg P$ 与原命题等价，逆命题 $Q\Rightarrow P$ 一般不等价。

全称量词 $\forall$ 表示“对所有”，存在量词 $\exists$ 表示“至少有一个”。否定时既要否定谓词，也要交换量词：
\[
 \neg(\forall x\in D,\ P(x))\Longleftrightarrow
 \exists x\in D,\ \neg P(x),\qquad
 \neg(\exists x\in D,\ P(x))\Longleftrightarrow
 \forall x\in D,\ \neg P(x).
\]
“对每个实数 $x$，存在实数 $y>x$”为真，可取 $y=x+1$；交换量词成为“存在实数 $y$ 大于每个实数 $x$”则为假，取 $x=y$ 即违反要求。变量的依赖关系因此是命题的一部分。

\section{直接证明、反证与反例}
直接证明从假设与定义出发。例如偶数 $a=2m,b=2n$ 的和为 $2(m+n)$，仍为偶数；这里 $m,n\in\mathbb Z$ 保证 $m+n$ 是整数。证明“整数 $n$ 的平方为偶数则 $n$ 为偶数”，可证明逆否命题：若 $n=2k+1$，则 $n^2=2(2k^2+2k)+1$ 为奇数。

\keyterm{反证法}（proof by contradiction）假定结论不成立，并从这一假定与原条件推出矛盾。它不能偷偷改动原条件，也不能把待证结论作为推理依据。否定全称命题只需一个满足全部假设却违反结论的反例，例如 $x=-1$ 否定了“$x^2=1$ 必有 $x=1$”；验证许多个正例仍不能证明该命题对所有对象成立。

\section{数学归纳法与递归描述}
\begin{theorem}[数学归纳法]
设 $P(n)$ 是对整数 $n\ge n_0$ 定义的命题。若 $P(n_0)$ 成立，且对每个 $n\ge n_0$ 都有 $P(n)\Rightarrow P(n+1)$，则所有 $n\ge n_0$ 上的 $P(n)$ 均成立。
\end{theorem}
\begin{proof}
采用正整数的良序性质，即非空正整数集合有最小元素。若存在反例，取最小反例 $m$。起点已验证，故 $m>n_0$；于是 $P(m-1)$ 成立，归纳步骤又推出 $P(m)$，矛盾。对任意整数起点，可先平移指标到正整数。
\end{proof}
例如 $1+\cdots+n=n(n+1)/2$：$n=1$ 时成立；若对 $n$ 成立，则增加 $n+1$ 后得到 $(n+1)(n+2)/2$，完成归纳。归纳假设并未假定所有情形成立，只把某一步成立作为推出下一步的条件。强归纳法允许在推导 $P(n+1)$ 时使用从起点到 $n$ 的全部结论，证明同样依赖最小反例。

递归定义则规定对象如何逐步生成，如 $a_1=1,\ a_{n+1}=a_n+2$ 唯一确定正奇数列。初值不可省略，否则许多不同数列可能遵守同一更新规则。

\section{例题、习题与学习检查}
\begin{example}
设 $A=\{x\in\mathbb R:x^2=1\}$，$B=\{x\in\mathbb R:x>0\}$。求同时或恰有一个条件成立的实数。
\end{example}
\begin{solve}
$A=\{-1,1\}$，故同时成立的集合为 $A\cap B=\{1\}$；恰有一个成立的集合为 $\{-1\}\cup((0,+\infty)\setminus\{1\})$。这里 $1$ 虽属于并集，却应从“恰有一个”的答案中剔除。
\end{solve}
\begin{enumerate}
\item 否定命题“对每个正整数 $n$，存在正整数 $m>n$ 且 $m$ 为偶数”。答：存在正整数 $n$，使每个正整数 $m$ 都满足 $m\le n$ 或 $m$ 为奇数。
\item 证明 $1+3+\cdots+(2n-1)=n^2$。提示：归纳步骤增加 $2n+1$。
\item 判断 $A\subseteq B\cup C$ 是否必有 $A\subseteq B$ 或 $A\subseteq C$。反例为 $A=\{1,2\},B=\{1\},C=\{2\}$。
\end{enumerate}

\chapter{数系、运算与近似}
\label{chap:em-numbers}
数系的扩展来自运算和测量的需要：减法引入负数，除法引入分数，几何长度进一步要求无理数。本章把实数的四则运算、顺序和完备性作为基础结构，重点研究正确运算与可靠近似；实数的构造属于数学分析的基础问题。

\section{整数、分数与有理数}
\keyterm{有理数}（rational number）是可写成 $p/q$ 的数，其中 $p,q\in\mathbb Z$、$q\ne0$。分数的不同表示满足
\[
 \frac pq=\frac rs\Longleftrightarrow ps=rq,\qquad
 \frac pq+\frac rs=\frac{ps+rq}{qs},\qquad
 \frac pq\frac rs=\frac{pr}{qs}.
\]
约分改变表示而不改变数值，但只能约去共同的非零因子。除法是乘以倒数，除数为零时没有定义，因为 $0\cdot x=b$ 在 $b\ne0$ 时无解，在 $b=0$ 时也不能唯一确定 $x$。

比例比较必须统一单位。若某比例由 $20\%$ 增至 $25\%$，其增加量为 $5$ 个百分点，而相对增加为 $(25-20)/20=25\%$。若商品先涨价 $10\%$ 再降价 $10\%$，最终价格乘以 $1.1\times0.9=0.99$，比原价低 $1\%$，两次百分比作用的基数不同。

\section{无理数与实数轴}
不属于 $\mathbb Q$ 的实数称为\keyterm{无理数}（irrational number）。以 $\sqrt2$ 为例，若 $\sqrt2=p/q$ 且 $p,q$ 已约至互素，则 $p^2=2q^2$。平方为偶数蕴含本身为偶数，故 $p=2k$；代入得 $q^2=2k^2$，又推出 $q$ 为偶数，与互素矛盾。最简分数的存在可由正分母取最小值得到，不需要预先使用素因数分解。

有理数的小数表示有限或最终循环：长除法中的余数只有有限种，余数为零则终止，否则必有余数重复。反之，最终循环小数可通过乘以适当的 $10$ 的幂再相减转化为分数，例如 $x=0.\overline{27}$ 满足 $100x-x=27$，故 $x=3/11$。小数表示本身未必唯一，如 $0.999\ldots=1$；无限小数的值需要通过极限或实数完备性解释。

\keyterm{实数}（real number）与数轴上的点相对应，$\mathbb R=\mathbb Q\cup(\mathbb R\setminus\mathbb Q)$。本书采用实数的完备性：每个非空且有上界的实数集合有最小上界。它保证数轴没有有理数系中那样的空缺；该性质并非仅由四则运算律推出。

\section{顺序、绝对值与距离}
实数的顺序与加法相容；乘以正数保持大小关系，乘以负数反向。定义
\[
 |x|=\begin{cases}x,&x\ge0,\\-x,&x<0.\end{cases}
\]
于是 $|x|\ge0$、$|xy|=|x||y|$，且 $-|x|\le x\le|x|$。两数 $x,y$ 在数轴上的距离为 $|x-y|$。对 $r>0$，$|x-a|<r$ 等价于 $a-r<x<a+r$，$|x-a|\ge r$ 等价于 $x\le a-r$ 或 $x\ge a+r$。

\begin{theorem}[三角不等式与反三角不等式]
对任意实数 $x,y$，
\[
 |x+y|\le |x|+|y|,\qquad
 \bigl||x|-|y|\bigr|\le |x-y|.
\]
第一式等号成立当且仅当 $xy\ge0$。
\end{theorem}
\begin{proof}
将 $-|x|\le x\le|x|$ 与对应的 $y$ 不等式相加，得到第一式。两侧非负，其平方之差为 $2(|xy|-xy)$，故等号条件为 $xy\ge0$。再由 $x=(x-y)+y$ 得 $|x|-|y|\le|x-y|$，交换 $x,y$ 即得第二式。
\end{proof}

\section{整数幂、根式与运算范围}
正整数幂由重复相乘定义；当 $a\ne0$ 时，规定 $a^0=1$、$a^{-n}=1/a^n$。指数律 $a^ma^n=a^{m+n}$ 与 $(a^m)^n=a^{mn}$ 在相应运算有定义时成立。本书不把 $0^0$ 作为通常实数幂定义。

当 $a\ge0$ 时，$\sqrt a$ 是平方等于 $a$ 的唯一非负实数，因此 $\sqrt{x^2}=|x|$。偶次实根要求被开方数非负并取非负值，奇次实根对所有实数存在且唯一。对 $a,b\ge0$ 有 $\sqrt{ab}=\sqrt a\sqrt b$，因为右侧非负且平方为 $ab$；若不满足范围条件，不能照搬公式。

有理化利用平方差，例如 $a,b\ge0$ 且 $a\ne b$ 时，
\[
 \frac1{\sqrt a-\sqrt b}=\frac{\sqrt a+\sqrt b}{a-b}.
\]
它既可用于化简，也能避免两个接近的近似数相减所造成的相对误差放大。

\section{近似数、有效数字与误差}
设准确值为 $x$，近似值为 $\widetilde x$。绝对误差为 $|\widetilde x-x|$，在 $x\ne0$ 时相对误差为 $|\widetilde x-x|/|x|$。保留小数点后 $m$ 位并四舍五入时，绝对误差至多为 $\frac12\,10^{-m}$；有效数字则从首个非零数字起计，反映相对精度的数量级。

若 $|x-\widetilde x|\le\varepsilon$、$|y-\widetilde y|\le\delta$，则三角不等式给出
\[
 |(x+y)-(\widetilde x+\widetilde y)|\le\varepsilon+\delta,
 \quad
 |xy-\widetilde x\widetilde y|
 \le|\widetilde x|\delta+|\widetilde y|\varepsilon+\varepsilon\delta.
\]
第二式由 $x=\widetilde x+u,y=\widetilde y+v$ 展开得到。商的误差控制还要求分母远离零；仅知道分母误差很小并不足够。计算通常保留若干保护位，在最终结果处统一舍入。

\section{例题、习题与学习检查}
\begin{example}
不使用小数近似，化简 $\sqrt{(x-2)^2}$；再估计 $\sqrt2$ 到三位小数。
\end{example}
\begin{solve}
第一式等于 $|x-2|$，即 $x\ge2$ 时为 $x-2$，$x<2$ 时为 $2-x$。又 $1.414^2=1.999396<2$，$1.415^2=2.002225>2$，故 $1.414<\sqrt2<1.415$。进一步有 $1.4145^2=2.00081025>2$，所以四舍五入到三位小数为 $1.414$，仅有前两个界还不能确定舍入方向。
\end{solve}
\begin{enumerate}
\item 化 $0.1\overline6$ 为分数。答：$1/6$。
\item 举例说明两无理数之和、之积都可能为有理数。可取 $\sqrt2,-\sqrt2$，其和为 $0$、积为 $-2$。
\item 长方形两边近似为 $3,4$，各有绝对误差不超过 $0.01$，求面积误差界。答：相对于近似面积 $12$，误差不超过 $0.0701$。
\end{enumerate}

\chapter{整除、素数与同余初步}
\label{chap:em-integers}
整数问题同时受到运算关系和取值离散性的约束。方程在实数中可解，并不保证在整数中可解；整除、余数与最大公因数正是表达这类约束的基本工具。

\section{整除与带余除法}
对整数 $a,b$，若存在 $k\in\mathbb Z$ 使 $b=ak$，称 $a$ \keyterm{整除}（divide）$b$，记为 $a\mid b$。若 $d\mid a,d\mid b$，则对任意整数 $u,v$ 有 $d\mid(ua+vb)$，因为整数倍之和仍为整数倍。

\begin{theorem}[带余除法]
对整数 $a$ 和正整数 $d$，存在唯一整数 $q,r$ 使 $a=qd+r$ 且 $0\le r<d$。
\end{theorem}
\begin{proof}
集合 $\{a-kd:k\in\mathbb Z,\ a-kd\ge0\}$ 非空，取其最小元素 $r=a-qd$。若 $r\ge d$，则 $r-d$ 也是其中更小的元素，矛盾。若另有 $a=q'd+r'$ 且 $0\le r'<d$，则 $(q-q')d=r'-r$，右端绝对值小于 $d$，只能为零。
\end{proof}
例如 $-17=(-4)\cdot5+3$。负数除法也按非负余数约定，商不是把 $-17/5$ 向零截断所得的数。

\section{最大公因数与欧几里得算法}
不全为零的整数 $a,b$ 的最大正公因数记为 $\gcd(a,b)$；若其为 $1$，称两数互素。由于 $a,b$ 与 $b,a-qb$ 的公因数集合相同，
\[
 \gcd(a,b)=\gcd(b,a-qb).
\]
\keyterm{欧几里得算法}（Euclidean algorithm）以 $|a|,|b|$ 为输入，第二数非零时重复用带余除法得到的余数替换它，同时以前一第二数替换第一数；第二数为零时输出第一数。非负余数严格下降，保证算法有限步终止。

每个余数是前两个数的整数线性组合，反向代入得到 \keyterm{Bézout 等式}（Bézout identity）
\begin{equation}\label{eq:em-bezout}
 ua+vb=\gcd(a,b),\qquad u,v\in\mathbb Z.
\end{equation}
例如 $84=2\cdot30+24$、$30=24+6$，故 $\gcd(84,30)=6=3\cdot30-84$。算法同时计算公因数并提供线性组合的构造方法。

\section{素数与唯一分解}
大于 $1$ 且正因数只有 $1$ 与自身的整数称为\keyterm{素数}（prime number）；大于 $1$ 的其他整数称为合数。若素数 $p$ 不整除 $a$，则 $\gcd(p,a)=1$；由式 \eqref{eq:em-bezout}，存在 $u,v$ 使 $up+va=1$，乘以 $b$ 可知
\[
 p\mid ab\quad\Longrightarrow\quad p\mid a\ \text{或}\ p\mid b.
\]
\begin{theorem}[算术基本定理]
每个大于 $1$ 的整数都能写成素数乘积，且除因子的排列次序外表示唯一。
\end{theorem}
\begin{proof}
存在性用强归纳法：素数已符合要求，合数 $n=ab$ 中 $1<a,b<n$，分别分解 $a,b$ 即可。若同一数有两种素数分解，第一种的任一素因子整除第二种的乘积，反复使用上述性质可知它等于其中某一素因子。约去该因子继续比较，最终各因子及重数一致。
\end{proof}
素数有无穷多个：若全部素数为 $p_1,\ldots,p_k$，则 $p_1\cdots p_k+1>1$ 必有素因子，却不被其中任何一个整除，矛盾。若把正整数 $a,b$ 的素因数幂次分别列出，则最大公因数取各幂次的较小者，最小公倍数取较大者，所以 $\gcd(a,b)\operatorname{lcm}(a,b)=ab$。

\section{同余与余数计算}
对正整数 $m$，若 $m\mid(a-b)$，称 $a,b$ 模 $m$ \keyterm{同余}（congruent），记为 $a\equiv b\pmod m$。它等价于两数除以 $m$ 的余数相同。由整除的线性组合性质，同余可相加、相减、相乘，因而可取非负整数次幂。

同余不能任意约去因子。例如 $2\cdot1\equiv2\cdot3\pmod4$，却有 $1\not\equiv3\pmod4$。若 $\gcd(c,m)=1$，则由 Bézout 等式可把 $c$ 消去：$m\mid c(a-b)$ 蕴含 $m\mid(a-b)$。十进制整数模 $9$ 等于各位数字和模 $9$，因为 $10\equiv1\pmod9$；求末位则在模 $10$ 下计算幂。

\section{整数方程与整除论证}
\begin{theorem}
设 $a,b$ 为正整数，$c\in\mathbb Z$，$d=\gcd(a,b)$。方程 $ax+by=c$ 有整数解当且仅当 $d\mid c$。若 $(x_0,y_0)$ 是一个解，则全部整数解为
\[
 x=x_0+\frac bd t,\qquad y=y_0-\frac ad t,\qquad t\in\mathbb Z.
\]
\end{theorem}
\begin{proof}
必要性来自 $d$ 整除左端；充分性由式 \eqref{eq:em-bezout} 乘以整数 $c/d$ 得到。任意两解之差满足 $(a/d)(x-x_0)=-(b/d)(y-y_0)$。两系数互素，因此 $b/d$ 整除 $x-x_0$，得到上述表示；反向代入验证每个这样的数对都是解。
\end{proof}
若还要求非负解，只需对参数 $t$ 加上对应的两个线性不等式。模运算也能排除解：整数平方模 $4$ 只可能为 $0$ 或 $1$，故 $x^2+y^2=4k+3$ 没有整数解。

\section{例题、习题与学习检查}
\begin{example}
求 $6x+15y=21$ 的全部非负整数解。
\end{example}
\begin{solve}
除以 $3$ 得 $2x+5y=7$，特解为 $(1,1)$，全部整数解为 $x=1+5t,y=1-2t$。非负要求 $t\ge-1/5$ 且 $t\le1/2$，整数 $t$ 只能为 $0$，故答案为 $(1,1)$。
\end{solve}
\begin{enumerate}
\item 求 $\gcd(252,198)$ 及一个 Bézout 表示。答：$18=4\cdot252-5\cdot198$。
\item 求 $7^{2026}$ 的末位数字。答：$9$；幂的末位以 $7,9,3,1$ 循环。
\item 证明两个连续整数互素，并据此证明 $n(n+1)$ 为偶数。第一问取公因数整除两者之差，第二问也可直接按奇偶分类。
\end{enumerate}

\chapter{代数式与多项式}
\label{chap:em-expressions}
代数式把具体数量关系写成可重复使用的结构。一个可靠的变形既要保持数值关系，也要保留变量范围；因式分解、配方和换元的目的，是使表达式更适合当前的求解或证明任务。

\section{变量、代数式与恒等关系}
变量表示某个给定范围内可变的数，参数在一次求解中视为固定，但可用于比较一族问题。\keyterm{恒等式}（identity）是在所声明范围内对每个变量取值都成立的等式，例如 $(x+1)^2=x^2+2x+1$ 对一切实数成立；$x^2=1$ 则只对特定值成立。

表达式的定义域是使其全部运算合法的取值集合。分式要求分母非零，偶次根式要求被开方数非负。$\frac{x^2-1}{x-1}=x+1$ 在 $x\ne1$ 时成立；右端虽然在 $x=1$ 也有值，却不能据此扩大原式的定义域。这一区分在连续延拓与解方程时尤其重要。

\section{多项式的运算与恒等式}
实系数一元\keyterm{多项式}（polynomial）写为 $f(x)=a_0+a_1x+\cdots+a_nx^n$。若 $a_n\ne0$，则次数为 $n$、首项系数为 $a_n$；零多项式不赋予通常的非负整数次数。由分配律可得
\[
 (a\pm b)^2=a^2\pm2ab+b^2,\quad
 a^2-b^2=(a-b)(a+b),\quad
 a^3\pm b^3=(a\pm b)(a^2\mp ab+b^2).
\]
这些公式都是运算律的结果，而非额外假设。非零多项式的乘积次数等于次数之和，因为最高次项系数相乘非零；和的次数却可能因首项抵消而降低。

多项式作为形式表达式相等，要求每一幂次的系数相同。把它们视为实数函数时，在所有实数上取值相同也能推出系数相同，依据是本章后面证明的根数界；仅比较几个数值一般不足以认定恒等。

\section{因式分解的结构方法}
分解先寻找公共因子，再观察平方差、完全平方或可分组的结构。例如
\[
 x^3+x^2-x-1=(x+1)(x^2-1)=(x-1)(x+1)^2.
\]
配方把二次式转化为平方与常量的组合；换元则把反复出现的子表达式作为整体，例如令 $t=x^2$ 可把 $x^4-5x^2+4$ 分解为 $(t-1)(t-4)$，再回代得到 $(x-1)(x+1)(x-2)(x+2)$。

分解依赖系数范围。$x^2-2$ 在有理数范围内没有一次因子，否则有理根将等于 $\pm\sqrt2$；在实数范围却可写成 $(x-\sqrt2)(x+\sqrt2)$。$x^2+1$ 在实数中无根，不能分解成实系数一次因子的乘积；复数会改变这一结论。

\section{多项式除法、余式与因式}
\begin{theorem}[多项式带余除法]
对实系数多项式 $f,g$ 且 $g\ne0$，存在唯一多项式 $q,r$ 使 $f=qg+r$，其中 $r=0$ 或 $\deg r<\deg g$。
\end{theorem}
\begin{proof}
当被处理多项式次数不小于 $\deg g$ 时，减去适当单项式倍的 $g$ 消去最高次项。每次次数严格下降，有限步后得到所需余式。若有两个表示，则 $(q-q')g=r'-r$；若 $q-q'\ne0$，左端次数至少为 $\deg g$，右端却小于它，矛盾。
\end{proof}
取 $g(x)=x-a$，余式为常数，代入 $x=a$ 得余式为 $f(a)$。于是\keyterm{因式定理}（factor theorem）为
\begin{equation}\label{eq:em-factor}
 f(a)=0\Longleftrightarrow f(x)=(x-a)q(x).
\end{equation}
非零 $n$ 次多项式至多有 $n$ 个不同实根：$n=0$ 时没有根；若 $a$ 为一根，除去因子 $x-a$ 后，其余不同根都是次数 $n-1$ 的商的根，归纳完成。因此两个次数至多为 $n$ 的多项式，若在 $n+1$ 个不同点取值相同，其差必为零多项式。

\section{分式、根式与等价化简}
分式加减先通分，乘除先检查因子与除数是否为零。例如
\[
 \frac1{x-1}-\frac1{x+1}=\frac2{x^2-1},\qquad x\ne\pm1.
\]
约分只约乘积中的公共因子，不能把和中的一项随意删去。根式化简还需保留符号，如 $\sqrt{x^2y}=|x|\sqrt y$ 在 $y\ge0$ 时成立；直接写成 $x\sqrt y$ 会在 $x<0,y>0$ 时出错。

有理化应服从目的。表达式 $\sqrt{x+1}-\sqrt x$ 在 $x\ge0$ 时等于 $1/(\sqrt{x+1}+\sqrt x)$；后者既显示它为正，也便于比较大小。这里分母严格为正，因此转换没有丢失取值。

\section{例题、习题与学习检查}
\begin{example}
求 $f(x)=x^3-2x^2-x+2$ 的全部实根，并求 $f(x)$ 除以 $x-3$ 的余式。
\end{example}
\begin{solve}
分组得 $f(x)=x^2(x-2)-(x-2)=(x-2)(x-1)(x+1)$，故根为 $2,1,-1$。余式由代入得到 $f(3)=27-18-3+2=8$。求根与求余式使用不同结构，不必为了余式完整做一次除法。
\end{solve}
\begin{enumerate}
\item 若 $x^2+ax+b=(x+2)^2$ 恒成立，求 $a,b$。答：$a=4,b=4$。
\item 化简 $(x^2-4)/(x^2+x-6)$ 并保留条件。答：$(x+2)/(x+3)$，$x\ne2,-3$。
\item 一个次数至多为 $3$ 的多项式在 $0,1,2,3$ 处均取值 $5$，证明它恒等于 $5$。提示：研究它与常数 $5$ 的差。
\end{enumerate}

\part{方程、不等式与函数}

\chapter{方程与方程组}
\label{chap:em-equations}
方程求解的目标是确定全部解，而不是把式子变得好看。每一次操作都应回答两个问题：原解是否仍被保留，新式是否产生原来不允许的解。以下除特别说明外，未知数均在实数范围内取值。

\section{方程、解集与等价变形}
\keyterm{方程}（equation）是要求未知数满足的等式，解集由定义域内使等式成立的全部值组成。两方程解集相同才称等价。在公共定义域内，同加同减保持等价；同乘一个可能为零的式子只保证原解仍成立，同除该式子则可能丢解。

平方通常是单向操作：$u=v$ 推出 $u^2=v^2$，逆向只能得到 $u=v$ 或 $u=-v$。若已经知道 $u,v\ge0$，平方才等价。求解时可以使用必要条件得到有限个候选值，但必须代回原方程筛选；验根是恢复充分性的步骤，并非装饰性检查。

\section{一元一次方程与线性方程组}
对实参数 $a,b$，方程 $ax=b$ 在 $a\ne0$ 时唯一解为 $b/a$，在 $a=0,b\ne0$ 时无解，在 $a=b=0$ 时每个实数都是解。参数退化必须在除法前处理。

方程组的解要同时满足各方程。交换方程、将某一方程乘以非零常数、将一方程的倍数加到另一方程上，都是可逆操作，因而保持公共解集。对
\[
 ax+by=e,\qquad cx+dy=f,
\]
消元给出 $(ad-bc)x=ed-bf$。若 $ad-bc\ne0$，再回代可得
\[
 x=\frac{ed-bf}{ad-bc},\qquad y=\frac{af-ec}{ad-bc}.
\]
若 $ad-bc=0$，应继续检查方程是矛盾还是相关，而不能使用上述商式。一般线性方程组也按逐个消元、识别矛盾或自由变量、最后回代的顺序处理。

\section{一元二次方程及根与系数}
设 $a\ne0$。将 $ax^2+bx+c=0$ 乘以 $4a$ 并配方，得到 $(2ax+b)^2=b^2-4ac$。记\keyterm{判别式}（discriminant）为 $\Delta=b^2-4ac$，则
\begin{equation}\label{eq:em-quadratic}
 x=\frac{-b\pm\sqrt\Delta}{2a}\quad(\Delta\ge0).
\end{equation}
$\Delta>0$ 有两个不同实根，$\Delta=0$ 有一个二重根，$\Delta<0$ 没有实根。若两根按重数记为 $x_1,x_2$，展开 $a(x-x_1)(x-x_2)$ 并比较系数得
\[
 x_1+x_2=-\frac ba,\qquad x_1x_2=\frac ca.
\]
这就是二次方程的\keyterm{韦达关系}（Viète's formulas）。反之，和为 $s$、积为 $p$ 的两实数必须是 $t^2-st+p=0$ 的根，所以存在的条件是 $s^2-4p\ge0$。

\section{可化归的高次、分式与根式方程}
高次方程先观察可分解或换元的结构。例如 $x^4-5x^2+4=0$ 令 $t=x^2\ge0$，得 $(t-1)(t-4)=0$，故 $x=\pm1,\pm2$。新变量的范围与每个新值对应的原变量都不能省略。

分式方程先排除分母零点，在剩余定义域内乘以公分母才等价。例如 $(x^2-1)/(x-1)=0$ 的候选根为 $\pm1$，但 $1$ 不在定义域，故只有 $-1$。根式方程 $\sqrt{x+2}=x$ 要求 $x\ge0$，平方得 $(x-2)(x+1)=0$，只有 $x=2$ 合法。若在一开始没有记录右侧非负，就容易留下增根 $-1$。

\section{含参数方程与非线性方程组}
分类依据应是运算或解的性质发生变化的位置。例如 $(a-1)x=a^2-1$ 在 $a=1$ 时恒成立，在 $a\ne1$ 时唯一解为 $a+1$。若提前约去 $a-1$，就丢失了整个退化情形。

根的位置可以结合判别式与韦达关系判断。实系数二次方程的两实根均为正，当且仅当判别式非负、根和为正、根积为正；后两条件说明两根同号且不可能同时为负。非线性方程组也可利用对称关系：$x+y=5,xy=6$ 意味着 $x,y$ 是 $t^2-5t+6=0$ 的两根，因此解为 $(2,3),(3,2)$。这种方法先确定无序根，再恢复有序数对。

\section{例题、习题与学习检查}
\begin{example}
求使 $x^2-2mx+m+2=0$ 有两个不同正实根的实数 $m$。
\end{example}
\begin{solve}
不同实根要求 $\Delta=4(m-2)(m+1)>0$，即 $m<-1$ 或 $m>2$；根和为 $2m>0$，根积为 $m+2>0$。三条件取交集得 $m>2$。$m=2$ 虽给出正根 $2$，但为重根，不符合题设。
\end{solve}
\begin{enumerate}
\item 解 $\sqrt{2x+3}=x$。答：$x=3$。
\item 解 $x+y=7,\ x^2+y^2=25$。答：$(3,4),(4,3)$，先由平方和求积。
\item 判断 $x(x-1)=x$ 中两侧约去 $x$ 会漏掉什么。答：原解为 $0,2$，约去后漏掉 $0$。
\end{enumerate}

\chapter{不等式与不等式组}
\label{chap:em-solving-inequalities}
解不等式是找出满足指定大小关系的取值；证明一个不等式则是说明它在给定范围内始终成立。两者都以顺序性质为基础，但任务不同。本章侧重解集与参数范围，普遍估计方法见第\ref{chap:em-inequality-methods}章。

\section{不等式的性质与解集运算}
若 $a<b$，则 $a+c<b+c$；乘以 $c>0$ 保持方向，乘以 $c<0$ 反向。原因是 $b-a>0$，而 $c(b-a)$ 的符号由 $c$ 决定。未知表达式作乘除因子时必须先判断符号，尤其不能在分式不等式中不加区分地乘以分母。

不等式组取各解集的交，分类得到的候选区间取并。严格不等式不包含取等端点，非严格不等式允许取等，但定义域排除的点永远不能加入。例如 $x>1$ 且 $x\le3$ 的解集为 $(1,3]$；$x<-1$ 或 $x>1$ 的解集是两个区间之并。

\section{一次与二次不等式}
一次不等式 $ax>b$ 在 $a>0$ 时为 $x>b/a$，在 $a<0$ 时为 $x<b/a$，在 $a=0$ 时根据 $0>b$ 的真假分别得到全体实数或空集。

对 $a\ne0$ 的二次式 $q(x)=ax^2+bx+c$，若 $\Delta<0$，配方
\[
 q(x)=a\left[\left(x+\frac b{2a}\right)^2-\frac{\Delta}{4a^2}\right]
\]
表明它始终与 $a$ 同号。若 $\Delta=0$，除重根外同样与 $a$ 同号，重根处为零。若有两不同根 $\alpha<\beta$，则 $q(x)=a(x-\alpha)(x-\beta)$，根外两因子同号，根内异号。这给出全部二次不等式的符号规则，而不必只靠抛物线图像记忆。

\section{多项式与分式的符号分析}
将可分解表达式写成常数与一次因子幂的乘积，并将分子零点、分母零点按大小排列。每个开区间内各因子的符号不变；越过奇重根时符号改变，越过偶重根时不变。分母零点即使同时是分子零点，也须从原定义域排除。

\begin{example}
解 $\dfrac{(x-4)^2(x+2)}{x-3}\le0$。
\end{example}
\begin{solve}
临界点为 $-2,3,4$，其中 $3$ 不允许取。除 $x=4$ 外，平方因子为正，故只需判断 $(x+2)/(x-3)$，它在 $[-2,3)$ 上非正；$x=4$ 使原式为零，也应加入。因此解集为 $[-2,3)\cup\{4\}$。偶重根可以产生孤立解，不能只记录变号点。
\end{solve}

\section{绝对值、根式与混合条件}
对 $r\ge0$，$|u|\le r$ 等价于 $-r\le u\le r$；若 $r<0$ 则无解。$|u|\ge r$ 在 $r\le0$ 时恒成立，在 $r>0$ 时等价于 $u\le-r$ 或 $u\ge r$。多个绝对值以各内部表达式的零点分段，逐段去绝对值，最终仍与该段范围相交。

根式比较须先建立定义域。求解 $\sqrt{x+2}>x$ 时，$x\ge-2$；在 $-2\le x<0$ 上右端为负，原式自动成立；在 $x\ge0$ 上可等价平方，得到 $x+2>x^2$，即 $0\le x<2$。合并为 $[-2,2)$。若直接平方并保留整个平方后的区间，将错误排除部分负数。

\section{含参数不等式与恒成立问题}
“存在 $x$ 使不等式成立”和“对所有 $x$ 都成立”具有不同量词。对 $a>0$ 的二次式，$ax^2+bx+c\ge0$ 对全部实数成立，当且仅当 $\Delta\le0$：配方后最小值为 $-\Delta/(4a)$。若 $a=0$，必须另作一次式或常数式讨论；若 $a<0$，不可能在全体实数上非负。

在闭区间上判断恒成立，则看该区间最小值。向上开口的二次式若顶点落在区间内，取顶点值；否则函数值随离顶点的距离增加，最小值在离顶点较近的端点取得。例如 $x^2-2tx+1\ge0$ 对所有 $x\in[0,1]$ 成立，$t\le0$ 时最小值为 $1$；$0<t<1$ 时为 $1-t^2>0$；$t\ge1$ 时为 $2-2t$，故总条件为 $t\le1$。

\section{例题、习题与学习检查}
\begin{example}
解 $|x-1|+|x+1|\le4$。
\end{example}
\begin{solve}
在 $x\le-1$ 时左端为 $-2x$，给出 $-2\le x\le-1$；在 $-1\le x\le1$ 时左端恒为 $2$；在 $x\ge1$ 时左端为 $2x$，给出 $1\le x\le2$。并集为 $[-2,2]$。各分段边界重复计入并不影响集合结果。
\end{solve}
\begin{enumerate}
\item 解 $(x-1)/(x+2)>0$。答：$(-\infty,-2)\cup(1,+\infty)$。
\item 求使 $x^2+2x+k\ge0$ 对所有实数 $x$ 成立的 $k$。答：$k\ge1$。
\item 解 $\sqrt{x+1}\le x-1$。答：$[3,+\infty)$；先有 $x\ge1$，再平方。
\end{enumerate}

\chapter{函数与图像}
\label{chap:em-functions}
方程关注哪些输入产生指定输出，函数则研究全部输入与输出之间的关系。同一关系可以用式子、图像和语言表达；这些表示应彼此解释，而不能用不精确的图形观察替代定义。

\section{函数、定义域与值域}
\begin{definition}
\keyterm{函数}（function）$f:D\to Y$ 给每个 $x\in D$ 指定唯一元素 $f(x)\in Y$。$D$ 称为定义域，$Y$ 称为陪域，集合 $f(D)=\{f(x):x\in D\}$ 称为值域。
\end{definition}
函数不必有一个统一解析式。例如分段计费可由几段规则共同规定。相同解析式配上不同定义域也能产生不同函数：$x^2$ 在 $\mathbb R$ 上值域为 $[0,+\infty)$，在 $[1,2]$ 上值域为 $[1,4]$。定义域限制来自题目条件与运算合法性二者的交集。

若规定函数相等还须具有相同陪域，就应把陪域也视为函数数据；本书讨论实值函数时统一把陪域取为 $\mathbb R$，重点比较定义域和对应规则。求值域需要说明每个候选输出确实能由某个允许输入取得，而非只给出一个包含输出的区间。

\section{图像、坐标与基本函数}
函数图像是点集 $\{(x,f(x)):x\in D\}$。每条竖直直线至多与它相交一次，这是单值性的几何表达。一次函数 $kx+b$ 的图像是直线；$k=0$ 时为水平直线。二次函数配方为
\[
 ax^2+bx+c=a\left(x+\frac b{2a}\right)^2+c-\frac{b^2}{4a},\qquad a\ne0,
\]
由平方非负即可确定顶点、对称轴和最值，不必预先使用导数。

反比例函数 $k/x$ 在 $k\ne0$ 时排除 $x=0$，取值也不为零；其两支的位置由 $k$ 的符号确定。$|x|$ 由 $x$ 与 $-x$ 两段组成。作图时先确定定义域和关键点，再研究区间内性质，不能只把少数采样点用任意光滑曲线连接。

\section{单调性、奇偶性与有界性}
在区间 $I$ 上，若任取 $x_1<x_2$ 都有 $f(x_1)<f(x_2)$，称 $f$ 严格递增；用 $\le$ 替换 $<$ 得到非严格递增。递减相应定义。例如在 $[0,+\infty)$ 上，
\[
 x_2^2-x_1^2=(x_2-x_1)(x_2+x_1)>0,
\]
故平方函数严格递增；在负半轴上方向相反。$1/x$ 在正、负半轴上分别严格递减，却不在整个不连通定义域上严格递减，例如 $-1<1$ 而 $-1<1$ 的函数值顺序没有反转。

定义域关于原点对称时，$f(-x)=f(x)$ 称为偶函数，$f(-x)=-f(x)$ 称为奇函数。奇函数只有在 $0$ 属于定义域时才能推出 $f(0)=0$。若存在 $M$ 使所有 $f(x)\le M$，称 $M$ 为上界；只有某个输入实际取得这个上界且它不小于任何函数值时，才称它为最大值。$f(x)=x$ 在 $(0,1)$ 上有界，却没有最大值或最小值。

\section{复合函数、反函数与图像变换}
复合函数 $(f\circ g)(x)=f(g(x))$ 的定义域是 $\{x\in D_g:g(x)\in D_f\}$。例如 $f(u)=\sqrt u$、$g(x)=x^2-1$ 的复合定义域为 $(-\infty,-1]\cup[1,+\infty)$。次序通常不能交换。

若 $f$ 在 $D$ 上为单射，即相同输出只能来自相同输入，则对每个 $y\in f(D)$ 可以唯一恢复 $x$，从而定义\keyterm{反函数}（inverse function）$f^{-1}:f(D)\to D$。平方函数在整个实轴上不单射，限制在 $[0,+\infty)$ 后反函数为 $\sqrt x$；限制在 $(-\infty,0]$ 后则为 $-\sqrt x$。反函数的图像是原图像关于直线 $y=x$ 的反射，$f^{-1}$ 不是倒数 $1/f$。

图像变换可以从点对应推导：$(x,f(x))$ 在水平右移 $h$ 后成为 $(x+h,f(x))$，故新式为 $y=f(x-h)$；竖直平移得到 $f(x)+k$，纵向缩放得到 $Af(x)$。$|f(x)|$ 把横轴下方图像翻到上方，$f(|x|)$ 则把原右半图像向左镜像。

\section{零点、交点与函数观点下的求解}
$f(x)=g(x)$ 等价于公共定义域内 $f-g$ 的零点问题；$f(x)>g(x)$ 表示前者图像位于后者上方。严格单调函数至多有一个零点，因为两个不同输入不能得到同一输出，但单调性不保证零点存在，例如 $2^x>0$，其性质见下一章。

常用的零点存在依据是连续函数的介值性质：若 $f$ 在 $[a,b]$ 连续且 $f(a)f(b)<0$，则 $(a,b)$ 内有零点。连续性意为输入充分接近时输出也充分接近；其严格定义及介值定理证明属于分析基础，本书使用时明确列出条件。多项式连续，因此可用端点符号与单调性共同证明它在某区间恰有一根，图像只用于帮助定位。

\section{例题、习题与学习检查}
\begin{example}
求 $f(x)=x^2-4x+5$ 在 $[0,3]$ 上的值域。
\end{example}
\begin{solve}
$f(x)=(x-2)^2+1$。最小值 $1$ 在 $x=2$ 取得；区间中离 $2$ 最远的点为 $0$，最大值为 $5$。对每个 $y\in[1,5]$，取 $x=2-\sqrt{y-1}\in[0,2]$ 即有 $f(x)=y$，故值域恰为 $[1,5]$；这一构造补足了“所有中间值均能取到”的论证。
\end{solve}
\begin{enumerate}
\item 求 $f(x)=1/(x-2)$ 的定义域和值域。答：分别为 $\mathbb R\setminus\{2\}$ 与 $\mathbb R\setminus\{0\}$。
\item 求 $f(x)=3x-2$ 的反函数。答：$f^{-1}(x)=(x+2)/3$。
\item 证明 $x^3+x=1$ 恰有一个实根。提示：作差证明严格递增，再在 $[0,1]$ 上使用介值性质。
\end{enumerate}

\chapter{幂、指数与对数}
\label{chap:em-exponential}
线性模型描述每一步增加固定数量，指数模型描述每一步乘以固定比例。对数则反过来回答“需要多少次比例变化才能达到目标”。二者互逆，是处理尺度和增长的基本工具。

\section{有理指数与实指数的扩展}
对 $a>0$、整数 $p$ 和正整数 $q$，定义 $a^{p/q}=\sqrt[q]{a^p}$。若 $p/q=r/s$，把两个正数的幂提升到 $qs$ 次，得到同样的 $a^{ps}=a^{rq}$，由正数幂的单射性知定义不依赖分数表示。利用同样方法可证明有理指数的加法与乘法法则。

对 $a>1$ 和实数 $x$，可以利用完备性定义
\[
 a^x=\sup\{a^r:r\in\mathbb Q,\ r\le x\}.
\]
集合非空且有上界：选择整数 $m\le x<M$，则 $a^m$ 是其中元素，$a^M$ 是上界。对 $0<a<1$ 定义 $a^x=(1/a)^{-x}$，对 $a=1$ 定义 $1^x=1$。这一构造与有理指数相容，并给出连续、保持指数律的唯一延拓；完整证明还需实数逼近理论，本章采用该延拓性质，不以有限位小数试算代替证明。底数为正是这些实指数法则的共同前提。

\section{幂函数与定义域}
\keyterm{幂函数}（power function）把指数固定，如 $x^2,x^{-1},x^{1/2}$。整数正幂在实轴上都有定义，负整数幂排除零；偶次根限制输入非负，奇次根可接受负数。一般实指数的 $x^\alpha$ 默认在 $x>0$ 上讨论，再按具体指数判断能否扩展到零或负数。

对 $x>0$，$\alpha>0$ 时 $x^\alpha$ 严格递增，$\alpha<0$ 时严格递减，$\alpha=0$ 时恒为 $1$。整数幂可以用因式分解直接证明，例如 $0<u<v$ 时 $v^n-u^n=(v-u)(v^{n-1}+\cdots+u^{n-1})>0$；实指数情形依赖上一节的延拓性质。实负底数的分数幂则须单独定义，不能任意交换幂次顺序，例如 $((-1)^2)^{1/2}=1$ 而 $(-1)^{2\cdot(1/2)}=-1$。

\section{指数函数与增长模型}
固定 $a>0,a\ne1$ 时，$f(x)=a^x$ 称为\keyterm{指数函数}（exponential function）。它的定义域为 $\mathbb R$，值域为 $(0,+\infty)$，经过 $(0,1)$；$a>1$ 时严格递增，$0<a<1$ 时严格递减，并满足
\[
 a^{x+y}=a^xa^y,\qquad a^{x-y}=\frac{a^x}{a^y}.
\]
底数大小决定相同时间间隔的比例变化。例如初值 $V_0>0$、每期增长率 $r>-1$ 恒定时，离散模型为 $V_n=V_0(1+r)^n$。把整数期数延拓为实变量是额外建模选择，并不意味着实际过程在每一瞬间都遵循同一曲线。

\section{对数、换底与数量级}
\begin{definition}
设 $a>0,a\ne1,x>0$。满足 $a^y=x$ 的唯一实数 $y$ 称为 $x$ 以 $a$ 为底的\keyterm{对数}（logarithm），记为 $\log_a x$。
\end{definition}
指数函数连续、严格单调且覆盖全部正数，保证定义中的存在与唯一。于是 $\log_a1=0,\log_a a=1$，以及
\begin{equation}\label{eq:em-loglaws}
 \log_a(xy)=\log_a x+\log_a y,\quad
 \log_a(x/y)=\log_a x-\log_a y,\quad
 \log_a(x^t)=t\log_a x
\end{equation}
对 $x,y>0,t\in\mathbb R$ 成立。例如令 $u=\log_a x,v=\log_a y$，则 $xy=a^{u+v}$，由唯一性得到第一式，其他两式同理。

若 $b>0,b\ne1$，则 $a^{\log_a x}=x$ 两侧取 $\log_b$ 给出换底公式 $\log_a x=\log_b x/\log_b a$。底数 $10$ 的对数反映数量级；底数为常数 $e$ 的对数记为 $\ln x$，$e$ 的一种构造见第\ref{chap:em-modeling}章。对数不把加法变为加法，例如 $\log_{10}(1+1)>0$，而 $\log_{10}1+\log_{10}1=0$。

\section{指数与对数方程、不等式}
指数函数单射，故 $a^{u(x)}=a^{v(x)}$ 等价于 $u(x)=v(x)$，前提是固定底数合法且两指数有定义。指数不等式的方向由底数决定；对数方程则先限制真数为正。

例如 $\log_2(x-1)+\log_2(x+1)=3$ 的定义域为 $x>1$，由式 \eqref{eq:em-loglaws} 化为 $x^2-1=8$，两候选值 $\pm3$ 中只留下 $3$。又 $\log_{1/2}(x-1)\ge2$ 先要求 $x>1$；由于底数小于 $1$，得到 $0<x-1\le1/4$，故 $1<x\le5/4$。忽略单调方向会得到完全相反的区间。

\section{例题、习题与学习检查}
\begin{example}
一笔自拟资金按每年 $5\%$ 的固定比例增长，至少经过多少个完整年度才翻倍？
\end{example}
\begin{solve}
要求 $1.05^n\ge2$，$n$ 为非负整数，所以 $n\ge\log_{1.05}2$。直接计算可验证 $1.05^{14}<2<1.05^{15}$，故最小完整年数为 $15$。该答案依赖增长率恒定且收益持续计入本金的假设，不是对实际收益的预测。
\end{solve}
\begin{enumerate}
\item 解 $3^{2x-1}=27$。答：$x=2$。
\item 解 $\log_2(x+2)-\log_2x=1$。答：$x=2$，定义域 $x>0$。
\item 解 $(1/3)^{x-1}<9$。答：$x>-1$。
\end{enumerate}

\part{数列、估计与有限方法}

\chapter{数列、递推与求和}
\label{chap:em-sequences}
数列把函数的输入限制为整数指标，适合描述逐次累积、递推计算与离散变化。求通项和求和的共同思路，是寻找相邻项之间能够消去或简化的关系。

\section{数列、通项与递推关系}
\keyterm{数列}（sequence）是定义在正整数或其初始有限段上的函数，第 $n$ 项记为 $a_n$。前 $n$ 项和记为 $S_n=\sum_{k=1}^n a_k$，于是 $a_1=S_1$，而 $n\ge2$ 时 $a_n=S_n-S_{n-1}$。后一个关系不能在未定义 $S_0$ 时直接用于 $n=1$；若约定空和 $S_0=0$，则可统一写出。

递推关系用先前项决定后续项，必须配有足够初值。例如 $a_{n+1}=2a_n$ 配上不同 $a_1$ 就产生不同数列。有限个已知项也不能唯一决定无限数列：在某个候选通项 $f(n)$ 上增加 $c(n-1)\cdots(n-m)$，前 $m$ 项不变，后续通常改变。因此发现规律之后，还要从题设条件证明它。

\section{等差数列与等比数列}
若 $a_{n+1}-a_n=d$，则累加相邻差得 $a_n=a_1+(n-1)d$。将正序与倒序的和相加，每对均为 $a_1+a_n$，因此
\[
 S_n=\frac{n(a_1+a_n)}2=na_1+\frac{n(n-1)}2d.
\]
若 $a_{n+1}=qa_n$，且 $q\ne0$，则 $a_n=a_1q^{n-1}$。对任意实数 $q\ne1$，错位相减得到
\begin{equation}\label{eq:em-geometric-sum}
 (1-q)S_n=a_1(1-q^n),\qquad
 S_n=\frac{a_1(1-q^n)}{1-q}.
\end{equation}
$q=1$ 时 $S_n=na_1$；$q=0$ 时从第二项起全为零，$S_n=a_1$，直接由递推解释即可。零项情形不宜用 $a_{n+1}/a_n$ 定义公比。

\section{裂项、错位相减与分组求和}
\keyterm{裂项求和}（telescoping sum）利用 $a_k=b_k-b_{k+1}$ 使中间项抵消，例如
\[
 \sum_{k=1}^n\frac1{k(k+1)}
 =\sum_{k=1}^n\left(\frac1k-\frac1{k+1}\right)
 =1-\frac1{n+1}.
\]
间隔不是 $1$ 时须保留相应数量的首尾项，例如
\[
 \sum_{k=1}^n\frac1{k(k+2)}
 =\frac12\left(1+\frac12-\frac1{n+1}-\frac1{n+2}\right).
\]
根式也可裂项：$1/(\sqrt{k+1}+\sqrt k)=\sqrt{k+1}-\sqrt k$，前 $n$ 项和为 $\sqrt{n+1}-1$。

若 $T_n=\sum_{k=1}^n kq^{k-1}$，$q\ne0,1$，则 $T_n-qT_n=1+q+\cdots+q^{n-1}-nq^n$，结合式 \eqref{eq:em-geometric-sum} 得
\[
 T_n=\frac{1-(n+1)q^n+nq^{n+1}}{(1-q)^2}.
\]
这说明错位相减不仅能求等比和，还能降低等差因子的次数；$q=1$ 改用等差和，$q=0$ 则直接按原来的多项式求和约定计算。

\section{简单递推的化归}
对 $a_{n+1}=pa_n+q$，若 $p=1$，则为公差 $q$ 的等差数列；若 $p\ne1$，令 $L=q/(1-p)$，有 $a_{n+1}-L=p(a_n-L)$。当 $p\ne0,1$ 时
\[
 a_n=L+p^{n-1}(a_1-L).
\]
$p=0$ 时从第二项起恒为 $q$，不必引入 $0^0$。这里 $L$ 是更新规则的不动点，平移把非齐次递推化为等比递推。

二阶例子 $F_1=F_2=1,F_{n+2}=F_{n+1}+F_n$ 的候选指数解满足 $r^2=r+1$。记 $\alpha=(1+\sqrt5)/2,\beta=(1-\sqrt5)/2$，直接代入初值与递推可验证 $F_n=(\alpha^n-\beta^n)/\sqrt5$。唯一性由递推逐项确定保证，因此验证完整初值和递推已经足以证明公式。

\section{单调、有界与极限的预备认识}
若对每个 $n$ 有 $a_{n+1}\ge a_n$，数列非严格递增；若还有统一上界，就能用完备性证明它收敛。这里 $a_n\to L$ 的严格含义是：对每个 $\varepsilon>0$，存在正整数 $N$，使所有 $n\ge N$ 都满足 $|a_n-L|<\varepsilon$。例如 $1/n\to0$，因为取整数 $N>1/\varepsilon$ 即可。

\begin{theorem}[单调有界收敛]
递增且有上界的实数列收敛到其所有项的上确界；递减且有下界的情形相应成立。
\end{theorem}
\begin{proof}
设 $L=\sup\{a_n:n\ge1\}$。对 $\varepsilon>0$，$L-\varepsilon$ 不是上界，故存在 $N$ 使 $a_N>L-\varepsilon$。递增性给出 $L-\varepsilon<a_N\le a_n\le L$，对所有 $n\ge N$ 成立，这就是极限定义。递减情形对 $-a_n$ 应用已证结果。
\end{proof}
当 $|q|<1$ 时 $q^n\to0$：$q=0$ 直接成立；否则令 $u=1/|q|-1>0$，归纳得 $(1+u)^n\ge1+nu$，所以 $|q|^n\le1/(1+nu)\to0$。因此等比级数的部分和在该条件下趋于 $a_1/(1-q)$；没有收敛条件，有限和公式不能直接变成无穷和。

\section{例题、习题与学习检查}
\begin{example}
设 $a_1=1,a_{n+1}=2a_n+3$，求通项。
\end{example}
\begin{solve}
不动点为 $L=-3$，故 $a_{n+1}+3=2(a_n+3)$，于是 $a_n+3=4\cdot2^{n-1}$，即 $a_n=2^{n+1}-3$。代入 $n=1$ 及更新式均吻合，确认初值和指标没有错位。
\end{solve}
\begin{enumerate}
\item 求 $\sum_{k=1}^n(2k-1)$ 与 $\sum_{k=1}^n3^{k-1}$。答：$n^2$ 与 $(3^n-1)/2$。
\item 设 $a_1=0,a_{n+1}=(a_n+2)/2$，求通项和极限。答：$a_n=2-2^{2-n}$，极限为 $2$。
\item 证明 $a_n=1-1/n$ 严格递增且极限为 $1$，但任何一项都不等于极限。
\end{enumerate}

\chapter{初等不等式与最值}
\label{chap:em-inequality-methods}
不等式把精确求值转化为范围控制，最值问题则进一步要求找到最好的界。有效证明必须同时交代变量范围、每步估计依据与取等条件，否则即使得到了正确的界，也可能没有解决最值问题。

\section{作差、作商与平方非负}
证明 $A\ge B$ 可以化为 $A-B\ge0$；若 $B>0$，也可以证明 $A/B\ge1$。最基本的非负结构是平方，例如
\[
 a^2+b^2-2ab=(a-b)^2\ge0,
\]
等号当且仅当 $a=b$。含参数的二次式常以配方处理，几个变量的式子则尝试写成若干平方的和。例如
\[
 a^2+b^2+c^2-ab-bc-ca
 =\tfrac12\bigl((a-b)^2+(b-c)^2+(c-a)^2\bigr)\ge0,
\]
等号当且仅当 $a=b=c$。给不同项分别估计之后，必须检查这些等号是否能够同时成立。

\section{常见平均数与均值不等式}
对正数 $x_1,\ldots,x_n$，定义调和平均（harmonic mean）、几何平均（geometric mean）、算术平均（arithmetic mean）与平方平均（quadratic mean）：
\[
 H=\frac n{\sum_i1/x_i},\quad G=\left(\prod_i x_i\right)^{1/n},
 \quad A=\frac1n\sum_i x_i,\quad Q=\sqrt{\frac1n\sum_i x_i^2}.
\]
\begin{theorem}[均值链]
$H\le G\le A\le Q$，任一相邻等号成立当且仅当所有 $x_i$ 相等。
\end{theorem}
\begin{proof}
二元关系 $\sqrt{uv}\le(u+v)/2$ 来自 $(\sqrt u-\sqrt v)^2\ge0$。反复两两分组可证明个数为 $2^k$ 时几何平均不超过算术平均。对一般 $n$，选择 $m=2^k\ge n$，在原数后补入 $m-n$ 个 $A$，新算术平均仍为 $A$，所以 $(\prod_i x_i)A^{m-n}\le A^m$，除以正数得到 $G\le A$。反复分组的等号要求全部数相等，给出取等条件。对倒数应用同一结论得 $1/G\le1/H$，即 $H\le G$。最后 $\sum_i(x_i-A)^2=\sum_i x_i^2-nA^2\ge0$，故 $A\le Q$，等号也要求每项等于 $A$。
\end{proof}
$G\le A$ 可推广到非负数，含零时直接判断即可；$H$ 仍需所有数为正。相同距离分别以速度 $u,v>0$ 行驶，平均速度为 $2/(1/u+1/v)$；相同时间行驶则为 $(u+v)/2$。选择平均数取决于总量如何累积。

\section{柯西--施瓦茨不等式与分式估计}
\begin{theorem}[柯西--施瓦茨不等式]
对实数 $a_i,b_i$（$1\le i\le n$），
\begin{equation}\label{eq:em-cauchy}
 \left(\sum_i a_ib_i\right)^2\le
 \left(\sum_i a_i^2\right)\left(\sum_i b_i^2\right).
\end{equation}
等号当且仅当对任意 $i,j$ 都有 $a_ib_j=a_jb_i$；若 $(b_i)$ 不全为零，等价于存在实数 $\lambda$ 使 $a_i=\lambda b_i$。
\end{theorem}
\begin{proof}
直接展开得到左、右两端之差的非负表达：
\[
 \left(\sum_i a_i^2\right)\left(\sum_i b_i^2\right)
 -\left(\sum_i a_ib_i\right)^2
 =\sum_{i<j}(a_ib_j-a_jb_i)^2.
\]
各平方同时为零即为等号条件。若某个 $b_j\ne0$，取 $\lambda=a_j/b_j$ 即得比例形式；若全部 $b_i=0$，两端同为零。
\end{proof}
对 $y_i>0$，代入 $a_i=x_i/\sqrt{y_i},b_i=\sqrt{y_i}$ 得
\[
 \sum_i\frac{x_i^2}{y_i}\ge\frac{(\sum_i x_i)^2}{\sum_i y_i},
\]
等号当且仅当各 $x_i/y_i$ 相同。该形式适合分母不同但总和可控的分式估计。

\section{排序、伯努利与凸性选读}
\keyterm{排序不等式}（rearrangement inequality）断言：若 $a_1\le\cdots\le a_n$、$b_1\le\cdots\le b_n$，则对任意排列 $\sigma$，
\[
 \sum_i a_ib_{n+1-i}\le\sum_i a_ib_{\sigma(i)}\le\sum_i a_ib_i.
\]
证明右侧只需消除逆序：交换 $i<j$ 上逆序的两个 $b$ 后，和的增加量为 $(a_j-a_i)(b_{\sigma(i)}-b_{\sigma(j)})\ge0$；有限次排序得到最大值，反向排列同理得到最小值。与最大值取等当且仅当不存在 $a_i<a_j$ 而所配 $b$ 值严格逆序的数对；两列均严格递增时，最大值只能由同序排列取得。

对整数 $n\ge1$、实数 $x\ge-1$，\keyterm{伯努利不等式}（Bernoulli's inequality）为 $(1+x)^n\ge1+nx$。归纳时乘以 $1+x\ge0$，再用
$(1+nx)(1+x)=1+(n+1)x+nx^2\ge1+(n+1)x$。$n=1$ 恒取等，$n>1$ 时仅 $x=0$ 取等。

区间 $I$ 上的函数 $f$ 称为\keyterm{凸函数}（convex function），若任意 $u,v\in I$、$0\le t\le1$ 都满足
$f((1-t)u+tv)\le(1-t)f(u)+tf(v)$。反复合并两个权重可归纳得到有限 Jensen 不等式
\[
 f\left(\sum_i\lambda_i x_i\right)\le\sum_i\lambda_i f(x_i),
 \quad \lambda_i\ge0,\quad\sum_i\lambda_i=1.
\]
归纳步骤先去掉零权重，再把前 $n-1$ 项归一化合为一个点，先后应用两点凸性与归纳假设。若两点不同且 $0<t<1$ 时严格不等，称严格凸；此时取等当且仅当正权重对应的 $x_i$ 全相等。$f(x)=x^2$ 的凸性直接由差值 $t(1-t)(u-v)^2$ 验证，无需微分法。

\section{约束最值与等号的可达性}
设矩形周长为 $L>0$，边长 $x,y>0$，则 $x+y=L/2$。均值不等式给出 $xy\le((x+y)/2)^2=L^2/16$；$x=y=L/4$ 满足约束且取等，所以它确为最大面积。

若约束改为 $x>1$，研究 $x+1/x$，虽有下界 $2$，但取等要求 $x=1$，不在可行域。取 $x=1+\varepsilon$ 并令正数 $\varepsilon$ 越来越小，可使函数值任意接近 $2$，因此 $2$ 是最大下界，却不是最小值。这解释了“得到一个界”和“证明最值”之间所差的可达性检验。

\section{例题、习题与学习检查}
\begin{example}
已知 $x,y>0$ 且 $x+y=1$，求 $1/x+1/y$ 的最小值。
\end{example}
\begin{solve}
由 $xy\le1/4$，有 $1/x+1/y=(x+y)/(xy)\ge4$；$x=y=1/2$ 时取等。也可直接使用分式柯西不等式得到 $1/x+1/y\ge(1+1)^2/(x+y)=4$，两种证明指向同一平衡条件。
\end{solve}
\begin{enumerate}
\item 若 $x,y\ge0,x+y=6$，求 $xy$ 最大值。答：$9$，在 $x=y=3$ 取得。
\item 证明 $(a+b+c)^2\le3(a^2+b^2+c^2)$ 并给出等号条件。答：应用式 \eqref{eq:em-cauchy}，等号为 $a=b=c$。
\item 用二点凸性的定义证明 $f(x)=|x|$ 为凸函数，并说明它不是严格凸函数。
\end{enumerate}

\chapter{计数原理与二项式定理}
\label{chap:em-counting}
计数的首要工作是准确规定对象：顺序是否重要、同一元素能否重复、哪些表示被视为同一个结果。公式只有在这些约定明确之后才有意义。

\section{加法原理、乘法原理与对应}
若有限集合 $A_1,\ldots,A_m$ 两两不交，则其并集大小为 $\sum_i|A_i|$，这称为加法原理。若一个选择过程分为 $m$ 步，且无论前面如何合法选择，第 $i$ 步都有恰好 $n_i$ 种选择，则总数为 $\prod_i n_i$，这称为乘法原理；若分支数不同，应分别计算再相加。

两个集合之间若能建立一一对应的\keyterm{双射}（bijection），则它们元素个数相同。例如从 $n$ 个对象中选出一个子集，等价于给每个对象指定选或不选，因此共有 $2^n$ 个子集。这里计数的是子集，而不是选择对象的先后过程。

\section{排列、组合与阶乘}
规定 $0!=1$，正整数 $n$ 的阶乘为 $n!=1\cdot2\cdots n$。从 $n$ 个不同对象中不重复地依次选 $k$ 个，$0\le k\le n$，有 $n(n-1)\cdots(n-k+1)=n!/(n-k)!$ 种。若不考虑顺序，每个 $k$ 元子集对应 $k!$ 个有序结果，所以\keyterm{组合数}（binomial coefficient）为
\[
 \binom nk=\frac{n!}{k!(n-k)!}.
\]
$\binom n0=1$ 对应唯一的空选择。

若 $n$ 个位置中有 $m_1,\ldots,m_r$ 个分别相同的对象，总数 $\sum_i m_i=n$，则不同排列数为 $n!/(m_1!\cdots m_r!)$；分母消去的是每个结果中相同的重复标记次数。$n\ge1$ 个不同对象围成圆环，若只把旋转视为相同而反射视为不同，固定一个对象为起点后有 $(n-1)!$ 种。改变等价关系会改变答案，不能再无条件除以 $2$。

\section{重复选择、分配与双阶乘}
非负整数解 $x_1+\cdots+x_k=n$（$n\ge0,k\ge1$）可用 $n$ 个星号与 $k-1$ 个隔板编码，隔板前后星号数量给出各 $x_i$，因而共有 $\binom{n+k-1}{k-1}$ 个。若要求每个 $x_i\ge1$，先令 $y_i=x_i-1$；当 $n\ge k$ 时共有 $\binom{n-1}{k-1}$ 个，否则无解。上界限制还需进一步筛选或容斥。

双阶乘每次递减 $2$ 相乘，约定 $0!!=1$。对 $n\ge1$，
\[
 (2n)!!=2^n n!,\qquad
 (2n-1)!!=\frac{(2n)!}{2^n n!}.
\]
将 $2n$ 个不同对象两两配对，固定第一个对象后它有 $2n-1$ 个搭档，再对剩余对象递归配对，故总数为 $(2n-1)!!$。也可从全排列中除以每对内部交换的 $2^n$ 与各对重新排序的 $n!$，得到同一公式。

\section{二项式定理与组合恒等式}
\begin{theorem}[二项式定理]
对非负整数 $n$，
\begin{equation}\label{eq:em-binomial}
 (x+y)^n=\sum_{k=0}^n\binom nk x^{n-k}y^k.
\end{equation}
\end{theorem}
\begin{proof}
当 $n\ge1$，展开 $n$ 个因子的乘积，每一项从每个因子中选 $x$ 或 $y$。恰选 $k$ 次 $y$ 的选择数为 $\binom nk$，于是得到对应系数。$n=0$ 时等式两端均为常数多项式 $1$，按多项式零次幂约定理解。
\end{proof}
把 $k$ 元子集按是否含指定元素分类，得到 Pascal 递推 $\binom nk=\binom{n-1}k+\binom{n-1}{k-1}$，边界之外的组合数约定为 $0$。取补集给出 $\binom nk=\binom n{n-k}$，式 \eqref{eq:em-binomial} 取 $x=y=1$ 得 $\sum_k\binom nk=2^n$；取 $x=1,y=-1$ 可比较偶数大小与奇数大小子集的数目。

\section{容斥与不重不漏的计数}
对有限集合 $A,B$，
\[
 |A\cup B|=|A|+|B|-|A\cap B|,
\]
因为交集中的元素在前两项被计了两次。三个集合相应满足
\[
 |A\cup B\cup C|=|A|+|B|+|C|
 -|A\cap B|-|A\cap C|-|B\cap C|+|A\cap B\cap C|.
\]
一般的\keyterm{容斥原理}（inclusion--exclusion principle）按交集大小交替加减。一个属于恰好 $r\ge1$ 个集合的元素得到的总系数为 $\sum_{j=1}^r(-1)^{j+1}\binom rj=1$，由二项式定理可知，因此每个并集元素恰计一次。

例如 $1$ 至 $100$ 中被 $2$ 或 $3$ 整除的整数有 $50+33-16=67$ 个，被两者整除等价于被 $6$ 整除。若要计算两者都不整除的数量，用全集大小减去并集即可，答案为 $33$。

\section{例题、习题与学习检查}
\begin{example}
把 $8$ 个相同物品分给 $3$ 个不同的人，允许有人分不到，且第一人至多分到 $3$ 个，共有多少种分法？
\end{example}
\begin{solve}
不限制上界时有 $\binom{10}2=45$ 种。违反限制意味着第一人至少有 $4$ 个，先给他取走 $4$ 个，余下 $4$ 个仍分给三人，共 $\binom62=15$ 种。故合法分法为 $45-15=30$ 种。相同物品没有内部排列，不能乘以 $8!$。
\end{solve}
\begin{enumerate}
\item 从 $6$ 人中选 $3$ 人并指定其中一人为负责人。答：$\binom63\cdot3=60$。
\item 将 $6$ 个不同对象两两配对。答：$5!!=15$。
\item 证明 $\sum_{k=0}^n k\binom nk=n2^{n-1}$（$n\ge1$）。提示：计数带有一个指定成员的非空子集，先选该成员再选其余成员。
\end{enumerate}

\chapter{有限问题中的证明方法}
\label{chap:em-discrete}
有限问题常常不需要复杂公式，而需要找到对象之间的组织方式。抽屉原理保证某种重复，不变量排除不可达状态，极端原理抓住最特殊对象，双重计数则把局部关系汇总为整体等式。

\section{抽屉原理与存在性}
\begin{theorem}[抽屉原理]
将 $N$ 个对象放入 $m\ge1$ 个盒子，至少一个盒子含有 $\lceil N/m\rceil$ 个或更多对象；$\lceil t\rceil$ 表示不小于 $t$ 的最小整数。
\end{theorem}
\begin{proof}
若每盒至多 $\lceil N/m\rceil-1$ 个，则总数至多为 $m(\lceil N/m\rceil-1)<N$，与对象总数矛盾。
\end{proof}
任取 $m+1$ 个整数，以模 $m$ 的余数作为盒子，则其中两数之差被 $m$ 整除。关键不是逐一计算差，而是识别只有 $m$ 种余数。存在性结论未必告诉我们具体是哪一对，若要找出对象，还需实际分组或搜索。

\section{不变量与单调量}
\keyterm{不变量}（invariant）是每次允许操作都保持的量。若初始与目标状态的不变量不同，目标就不可达。把棋盘按黑白交替染色，一块覆盖相邻两格的骨牌必覆盖一黑一白；从 $8\times8$ 棋盘删去同色的两个对角格后，黑白格数量相差 $2$，因此不能用骨牌无重叠覆盖。这一证明不必枚举所有铺法。

严格变化的量也能证明终止，但还需取值范围不能无限下降。欧几里得算法的余数是严格下降的非负整数，所以有限步终止；正实数 $1,1/2,1/4,\ldots$ 却能永远下降，故“每步减小”本身不足够。不变量一般只给必要条件：颜色数量相等不能单独保证任意形状都能铺满。

\section{极端原理与最小反例}
有限非空集合的实数值必能取得最大和最小值；正整数集合即使无限，只要非空也有最小元素。这使我们可以选择最小反例，再从它构造更小反例以得到矛盾，称为无限递降法。

\begin{example}
证明不存在正整数 $a,b,c$ 使 $a^2+b^2=3c^2$。
\end{example}
\begin{solve}
若存在，取 $c$ 最小的一组解。平方模 $3$ 只可能为 $0$ 或 $1$，两平方和模 $3$ 为零迫使 $a,b$ 都被 $3$ 整除，写成 $a=3u,b=3v$。代入得 $c^2=3(u^2+v^2)$，所以 $c=3w$。再约去 $9$ 得 $u^2+v^2=3w^2$，其中 $0<w=c/3<c$，与最小性矛盾。
\end{solve}
选择极端量时必须说明它存在，且构造后的对象仍满足全部条件；缺少后一检查，递降论证就不完整。

\section{双重计数与简单图论选读}
用两种方式数同一集合称为\keyterm{双重计数}（double counting）。有限无向简单图由顶点集与连接不同顶点的无向边组成，不含自环与重边；顶点的度是与它相接的边数。每条边恰有两个端点，所以
\[
 \sum_{v}\deg(v)=2|E|,
\]
其中 $E$ 为边集。左端奇偶性由奇度顶点数量决定，而右端为偶数，故奇度顶点个数必为偶数。

同一思想还能证明 $\binom n2=n(n-1)/2$：先数有序不同元素对得 $n(n-1)$，每个无序对被计两次。图的画法可以变化，边是否交叉及画得多长通常不属于图的组合数据；解释模型时应避免给图形增加未声明的几何条件。

\section{递推计数、构造与算法描述}
设 $T_n$ 为用 $1\times2$ 骨牌铺满 $2\times n$ 长方形的方法数。空区域有一种铺法，故 $T_0=1$，且 $T_1=1$。考察最左侧，或者放一块竖骨牌留下 $2\times(n-1)$ 区域，或者放两块横骨牌留下 $2\times(n-2)$ 区域，两类互斥且穷尽，因此
\[
 T_n=T_{n-1}+T_{n-2},\qquad n\ge2.
\]
依次维护相邻两个数即可算出 $T_n$，需要 $O(n)$ 次整数加法和常数个整数存储位置；若计入整数位数增长，则总位运算成本不是简单的 $O(n)$。

构造题除了证明目标不被不变量排除，还要给出合法操作序列，并说明何时结束。算法的正确性可以按循环步骤归纳，复杂度则需说明数的是运算次数、比较次数还是实际位操作。

\section{例题、习题与学习检查}
\begin{example}
单位正方形内任取五个点，证明其中两点的距离不超过 $\sqrt2/2$。
\end{example}
\begin{solve}
将正方形分成边长 $1/2$ 的四个小正方形，边界点按照固定规则归入其中一个。由抽屉原理有两点归入同一小方形，它们的距离不超过该方形对角线 $\sqrt{(1/2)^2+(1/2)^2}=\sqrt2/2$。
\end{solve}
\begin{enumerate}
\item 任取 $7$ 个整数，证明其中两数之差为 $6$ 的倍数。提示：按模 $6$ 余数分组。
\item 某图的顶点度数为 $1,1,2,2,4$，若该图存在，边数是多少？答：$5$；握手关系只给必要计数约束。
\item 计算 $T_5$ 并说明不能把两个第一步分支相乘。答：$T_5=8$，它们是互斥选择而非顺次步骤。
\end{enumerate}

\part{平面与空间几何}

\chapter{直线、三角形与相似}
\label{chap:em-triangles}
综合几何直接研究点、线、角和长度之间的关系。它的论证始于明确的几何基础，不始于图形“看上去如此”。本章采用欧氏平面的关联、顺序、长度和角度公理及通常的全等基础，不展开完整的公理化体系。

\section{基本对象、角与平行线}
两不同点确定唯一直线；线段包含两端点及其间的点，射线从端点沿一个方向延伸。长度非负并沿同一直线按次序可加，角的大小在指定范围内可加。平角为 $180^\circ$，两直线相交形成的对顶角相等，因为它们都等于平角减去同一邻角。

欧氏平行公设保证过直线外一点恰有一条平行于它的直线。平行线被第三条直线所截，同位角相等、内错角相等、同旁内角互补；反向判定也成立。角的相等关系必须来自这些条件或全等，而非来自线条画得近似平行。过三角形一个顶点作对边的平行线，即可把另外两个内角平移到该顶点附近，三角之和组成平角，故三角形内角和为 $180^\circ$。

\section{三角形的基本性质与全等}
\keyterm{全等}（congruence）表示两图形可由保持距离的移动完全重合。本书采用下列欧氏三角形基本判定：三边分别相等；两边及夹角分别相等；两角及其夹边分别相等。由内角和可推出“两角及任一对应边”判定，直角三角形还可由斜边与一条直角边判定。对应顶点的次序必须固定，夹角条件不能替换成任意角。

三角形中大边对大角、等边对等角，且任意两边之和大于第三边。后一性质也表达两点之间线段最短：绕经不在该线段上的第三点会增加距离。两边与一个非夹角不足以唯一确定三角形，可能出现两解；第\ref{chap:em-triangle-solving}章将给出其数量条件。退化的三点共线情形不属于本书所谓三角形。

\section{比例、相似与中位线}
两三角形对应角相等、对应边成同一正比例时称为\keyterm{相似}（similar）。常用判定为两角分别相等、两边成比例且夹角相等、三边成比例。其思路是把其中一个图形按对应边比例缩放到一对边相等，再归结到全等判定；缩放保持角度与长度比例是欧氏相似结构的基础性质。

设 $D\in AB,E\in AC$ 且 $DE\parallel BC$，对应角相等给出 $\triangle ADE\sim\triangle ABC$，所以 $AD/AB=AE/AC=DE/BC$。特别地，连接两边中点的中位线平行于第三边且长为其一半。相似比为 $k>0$ 时，底和高同时乘以 $k$，因而面积乘以 $k^2$。面积比不是一般的长度比，使用时须区分维度。

\section{勾股定理、面积与距离}
\begin{theorem}[勾股定理及逆定理]
直角三角形两直角边长为 $a,b$，斜边长为 $c$，则 $a^2+b^2=c^2$。反之，三角形三边满足此式时，长为 $c$ 的边所对的角为直角。
\end{theorem}
\begin{proof}
在直角顶点向斜边作高，垂足把斜边分成长 $p,q$ 的两段，$p+q=c$。两个小三角形分别与原三角形相似，得到 $a^2=cp,b^2=cq$，相加即得结论。反之，以长 $a,b$ 的两线段构造直角三角形，其斜边由已证公式等于 $c$，与原三角形三边相等，故原对应角也是直角。
\end{proof}
三角形面积为 $S=bh/2$，其中 $b$ 是所选底边长，$h$ 为对该底边所在直线的距离；这可由两份全等三角形拼成平行四边形得到。同高三角形面积比等于底边比，同底时等于高之比。面积法把多条线段关系转化为可加的量，往往比逐个求长更直接。

\section{特殊线段、交点与辅助线}
中线连接顶点与对边中点，角平分线把内角等分，高线垂直于对边所在直线。线段垂直平分线上的点到两端点等距，反向也成立：正向由直角三角形全等，反向由两等腰三角形或垂足两侧的勾股关系得到。

三边垂直平分线共点，因为前两条的交点到三个顶点距离相等，必也在第三条上，该点为外心。两内角平分线交点到三边距离相等，因而位于第三条内角平分线上，得到内心；角平分线等距性质由两个直角三角形全等证明。三条中线共点且把每条中线按从顶点起 $2:1$ 分割：两中线交点与两个边中点构成的三角形，利用中位线与原顶点边平行且为其一半，得到上述比；换另一对中线也得到同一唯一分点。该点称为重心。

三条高线也共点，称为垂心：分别过三个顶点作对边平行线形成外侧三角形，原三角形的顶点是外侧三角形各边中点，原高线就是外侧三角形的垂直平分线，归结到外心共点性。外心、垂心在钝角三角形中可位于外部。若 $AD$ 平分 $\angle A$ 且 $D\in BC$，则 $\triangle ABD,\triangle ACD$ 的面积比一方面为 $BD/DC$，另一方面以 $AB,AC$ 为底，由 $D$ 到两边距离相等得到 $AB/AC$，故 $BD/DC=AB/AC$。

\section{例题、习题与学习检查}
\begin{example}
直角三角形直角边为 $3,4$，求斜边上的高及它分出的两段长度。
\end{example}
\begin{solve}
斜边为 $5$，面积既为 $3\cdot4/2=6$，也为 $5h/2$，故 $h=12/5$。由相似关系，两投影长分别为 $3^2/5=9/5$ 和 $4^2/5=16/5$，其和为 $5$，可作为计算检查。
\end{solve}
\begin{enumerate}
\item 两相似三角形对应边比为 $2:3$，求周长比与面积比。答：$2:3$ 与 $4:9$。
\item 若三角形角平分线把对边分为 $2,3$，其相邻两边之和为 $10$，求两边。答：$4,6$。
\item 证明等腰三角形顶角平分线同时是底边上的中线与高线。提示：比较被分成的两个三角形。
\end{enumerate}

\chapter{四边形与圆}
\label{chap:em-circles}
四边形可以通过对角线分解为三角形，圆则借助半径、弦和圆周角连接不同三角形。选择辅助线时，应优先寻找已有的全等、相似和面积结构。

\section{四边形的性质与判定}
两组对边分别平行的四边形是平行四边形。作一条对角线，利用内错角与公共边可证两个三角形全等，进而得到对边相等、对角相等；比较两对角线交点附近的三角形，还可证明对角线互相平分。反之，若两对角线互相平分，由边角边全等可得对应内错角相等，从而两组对边平行。

矩形是有一个直角的平行四边形，菱形是四边相等的平行四边形，正方形同时是矩形与菱形。矩形对角线等长，菱形对角线互相垂直；结合平行四边形条件，各自的逆判定也成立。仅有对角线等长不足以保证矩形，等腰梯形就是反例。此处梯形约定为恰有一组对边平行的四边形。

\section{多边形的分割与面积}
从凸 $n$ 边形（$n\ge3$）的一个顶点向不相邻顶点连对角线，可得到 $n-2$ 个不重叠三角形，因此内角和为 $(n-2)180^\circ$。沿边界同向行走一周，转角总和为 $360^\circ$；对非凸图形须使用带方向的转角，不能把所有凹顶点都按正的小外角相加。

平行四边形面积等于底乘高，梯形两底长 $a,b$、高 $h$ 时面积为 $(a+b)h/2$，可由两份全等梯形拼成平行四边形证明。正多边形把中心与顶点相连，各小三角形具有共同的高 $\rho$，所以总面积为 $P\rho/2$，其中 $P$ 为周长，$\rho$ 为中心到边的距离。这个公式是圆面积逼近的重要桥梁。

\section{圆、弦与圆周角}
到定点 $O$ 距离等于 $r>0$ 的点集为圆，$O$ 为圆心、$r$ 为半径；圆及其内部构成圆盘。连接圆上两点的线段为弦，通过圆心的弦为直径。圆心到弦的垂线平分弦，因为垂足两侧的直角三角形斜边及一直角边相等；反向结论对非零弦也成立。

\begin{theorem}[圆周角定理]
圆上点 $P,A,B$ 两两不同，$\angle APB$ 的大小等于它所对的不含 $P$ 的弧 $AB$ 的圆心角大小的一半。
\end{theorem}
\begin{proof}
先考虑 $O$ 在角的一边上的情形，由两条半径组成的等腰三角形及外角定理得到所对圆心角为圆周角的两倍。一般情形连接 $PO$ 并延长为直径，将原角分为两个上述角或表示为它们之差；对应的弧也作同样的和或差，结论随之成立。这一处理同时覆盖圆心在角内和角外的情况。
\end{proof}
同弧上的圆周角相等，直径所对的圆周角为直角。圆内接四边形的两个对角所对弧合成整圆，故对角互补；反之，凸四边形一组对角互补可由圆周角的轨迹性质判定四点共圆。判定时须保留凸性与顶点次序。

\section{切线、割线与点的幂}
直线与圆恰有一个公共点时称为切线。切线垂直于过切点的半径：若半径不垂直，圆心到直线的垂足距离严格小于半径，直线就会有两个交点。反之，过圆上一点垂直于半径的直线也只有该交点，由勾股定理可见线上其他点到圆心更远。

设 $P$ 为圆外点，一条割线依次交圆于 $A,B$，另一条依次交于 $C,D$。由圆周角建立的相似三角形给出 $PA\cdot PB=PC\cdot PD$；若 $PT$ 为切线，则同样得到 $PT^2=PA\cdot PB$。过 $P,O$ 的割线进一步说明
\[
 PA\cdot PB=PT^2=OP^2-r^2.
\]
量 $OP^2-r^2$ 称为\keyterm{点的幂}（power of a point）。若 $P$ 在圆内且弦 $AB$ 经过它，无向长度则满足 $PA\cdot PB=r^2-OP^2$，其正值是点的幂的相反数，不能把圆外公式的符号直接搬入圆内。

\section{圆周长、面积与扇形}
考虑不断加倍边数的内接正多边形及对应外切正多边形。内周长增加，外周长减少，并提供共同上下界。若内接边长为 $\ell_n$、边数为 $n$、边心距为 $\rho_n$，则 $\rho_n^2=r^2-(\ell_n/2)^2$；内周长有界给出 $\ell_n\to0$，故 $\rho_n\to r$。外、内周长之比为 $r/\rho_n\to1$，因此两者趋于同一值 $C$，定义为圆周长。

相似缩放使 $C/r$ 与半径无关，记常数为 $2\pi$，得到 $C=2\pi r$。内、外面积分别为 $P_n\rho_n/2$ 与 $Q_nr/2$，都趋于 $Cr/2$；采用面积对这种穷竭逼近的相容性，圆盘面积为 $\pi r^2$。这说明公式依赖长度与面积的逼近理论，并非单凭拼接示意得出。圆心角为 $\alpha^\circ$、$0\le\alpha\le360$ 的扇形，由旋转对称、可加性及连续逼近，占整圆面积和弧长的比例均为 $\alpha/360$。

\section{例题、习题与学习检查}
\begin{example}
圆外点 $P$ 到圆心距离为 $13$，圆半径为 $5$，求切线段长；若一条割线上近交点距 $P$ 为 $8$，求远交点距离。
\end{example}
\begin{solve}
点的幂为 $13^2-5^2=144$，故切线段长为 $12$。割线关系给出 $8\cdot PB=144$，所以 $PB=18$；它大于近交点距离，符合交点次序。
\end{solve}
\begin{enumerate}
\item 菱形两对角线长为 $6,8$，求边长和面积。答：$5$ 与 $24$。
\item 半径 $3$ 的圆中，$120^\circ$ 扇形的弧长与面积是多少？答：$2\pi$ 与 $3\pi$。
\item 圆内两弦相交，其中一弦两段长为 $2,6$，另一弦一段长为 $3$，求另一段。答：$4$。
\end{enumerate}

\chapter{几何变换、轨迹与作图}
\label{chap:em-transformations}
几何变换通过移动或缩放整个图形，把难以直接比较的量放到同一位置。轨迹和作图则从相反方向出发：给定条件，寻找全部满足条件的点或构造出所需对象。

\section{平移、旋转与轴对称}
\keyterm{平移}（translation）把所有点沿同一方向移动同一距离；\keyterm{旋转}（rotation）固定中心 $O$，把各射线 $OP$ 转过同一有向角，并保持 $OP$ 长度；\keyterm{轴对称}（reflection）把点映到关于固定直线的镜像位置，连接点与像的线段被该直线垂直平分，轴上点保持不动。

这些变换都保持任意两点之间的距离，因而保持线段长度、角的大小、全等关系与面积。半周旋转就是中心对称。轴对称会颠倒平面定向，平移和旋转保持定向；因此“保持距离”不等于“保持顺逆时针次序”。变换复合依次执行，顺序一般不能交换，例如先旋转再沿固定方向平移通常不同于反过来操作。

\section{位似与相似变换}
给定中心 $O$ 和实数 $k\ne0$，\keyterm{位似}（homothety）把 $P$ 映为 $P'$：当 $k>0$ 时两点在同一射线 $OP$ 上且 $OP'=kOP$；当 $k<0$ 时在相反射线上且 $OP'=|k|OP$，中心仍映为自身。由相似三角形可得 $P'Q'=|k|PQ$，不经过中心的直线的像与原线平行。

正位似是缩放，负位似是缩放与中心对称的复合。与保持距离的变换复合后得到一般相似变换，它保持角的大小，把所有长度统一乘以同一正因子。圆的像仍为圆，半径乘以 $|k|$；面积乘以 $k^2$。比例 $k=0$ 会把所有点压到一点，不属于本节非退化相似变换。

\section{轨迹的必要性与充分性}
\keyterm{轨迹}（locus）是所有满足指定条件的点组成的集合。候选图形必须经过双向检验：每个合法点都在候选图形上，候选图形的每个允许点也满足原条件。只做到第一步，可能把真正轨迹扩大。

到定点 $O$ 距离为 $r>0$ 的点集是圆；到两个不同定点 $A,B$ 等距的点集是线段 $AB$ 的垂直平分线。到两条相交直线距离相等的点集是两条夹角平分线的并，而不是仅一条内角平分射线；若点被限制在某个角内部，才取其中对应分支。两不同平行线的等距点集则是它们之间的平行中线。

\section{尺规作图与存在条件}
尺规作图只允许使用无刻度直尺作两已知点连线，以及用圆规以已知点为圆心、已知长度为半径作圆，再取已作对象的交点。证明包括分析、构造、验证和解数讨论；能在图上估计一个点，不等于已给出尺规构造。

例如作 $AB$ 的垂直平分线，以 $A,B$ 为圆心取同一半径 $r>AB/2$ 作圆，交于 $P,Q$，连 $PQ$。两交点到 $A,B$ 等距，故均在垂直平分线上；两点确定直线，所作即目标。要求 $r>AB/2$ 是为了得到两个不同交点。

给定边长 $a,b,c>0$ 作三角形，可先作 $AB=c$，再取以 $A,B$ 为圆心、半径分别为 $b,a$ 的两圆交点。非退化交点存在当且仅当 $|a-b|<c<a+b$；两交点关于 $AB$ 对称，得到两种放置但同一个全等类型。等号情形共线，严格不等式失败则不可作成所需三角形。

\section{辅助构造、最短路径与方法比较}
\begin{theorem}[反射求最短折线]
设 $A,B$ 位于直线 $\ell$ 同一侧且都不在 $\ell$ 上。令 $A'$ 为 $A$ 关于 $\ell$ 的对称点，则 $P$ 遍历 $\ell$ 时，$AP+PB$ 的最小值为 $A'B$，在 $\ell$ 与线段 $A'B$ 的交点处取得。
\end{theorem}
\begin{proof}
对每个 $P\in\ell$，对称性给出 $AP=A'P$，所以 $AP+PB=A'P+PB\ge A'B$。$A',B$ 在直线两侧，连接线段恰穿过直线一次；在该交点处两段连接成一条直线且次序正确，三角不等式取等。
\end{proof}
若 $P$ 被进一步限制在直线的某个区段，交点可能不合法，须重新研究边界；不能只把无限直线上的答案照搬。辅助线的作用也是把原图形嵌入更容易使用全等、相似或三角不等式的结构，而不是无目的地增加线条。

\section{例题、习题与学习检查}
\begin{example}
在直角坐标图中，$A=(0,1),B=(4,3)$，折线 $APB$ 的转折点 $P$ 在横轴上。求最短总长。
\end{example}
\begin{solve}
把 $A$ 反射为 $A'=(0,-1)$，由勾股定理 $A'B=\sqrt{4^2+4^2}=4\sqrt2$。连接线的竖直变化从 $-1$ 到 $3$，与横轴相交时完成全程的四分之一，故 $P=(1,0)$，最短值确实可达。
\end{solve}
\begin{enumerate}
\item 位似比为 $-3$，长度和面积分别乘以多少？答：$3$ 与 $9$。
\item 给定边长 $2,3,6$ 能否作三角形？答：不能，因为 $2+3<6$。
\item 到两个不同点等距且到其中一个点距离为 $r$ 的点最多有几个？答：两个，由垂直平分线与圆的交点决定，也可能为零或一个。
\end{enumerate}

\chapter{立体几何与度量}
\label{chap:em-solid}
空间中的线面关系比平面复杂，尤其要区分相交与投影相交。解决度量问题的基本策略，是先确定空间位置，再截取包含目标长度或角的恰当平面，使用已经建立的平面几何。

\section{空间对象与相互位置}
本章采用欧氏空间的基本关联事实：不共线三点确定一个平面；直线上两点在某平面内，则整条直线在该平面内；两个不同平面若相交，其公共部分为一条直线。两不同直线可相交、平行或不共面，最后一种称为\keyterm{异面直线}（skew lines）。

直线与平面的关系为包含、相交一点或平行且无交点；两不同平面相交或平行。透视图中两线相交只表示投影相交，不能推出原空间直线有公共点。例如立方体一条底边与不在同一侧面上的竖棱可能为异面直线。

\section{平行与垂直的判定}
常用欧氏空间基本定理如下：平面外的一条直线若平行于平面内某一直线，则与该平面平行；一个平面内两条相交直线分别平行于另一平面内两条相交直线，则两平面平行。不能把“相交直线”减为两条平行直线，因为它们只提供一个方向。

直线垂直于平面，意为它垂直于经过垂足的每一条平面内直线；判定时只需检查其中两条相交直线。两平面垂直的常用充分条件，是其中一平面包含另一平面的一条垂线。这些判定作为欧氏空间的基础定理使用，完整的空间公理化证明不在本书展开；使用时必须确认有关直线确实属于指定平面。

例如长方体一个顶点出发的三条棱两两垂直，任意一条棱垂直于另外两棱所确定的面。已知只垂直于平面内一条直线却不够，因为该直线在平面内也有垂线。

\section{空间角、距离与截面}
异面直线夹角通过平移方向到共同点定义，通常取不超过 $90^\circ$ 的角。非垂直直线与平面的夹角是它与平面上正射影的夹角；垂直情形定义为 $90^\circ$。二面角的大小则在垂直于棱的截面中测量两半平面的夹角，范围为 $[0,180^\circ]$。

点到平面的距离为垂线段长度。若 $H$ 为垂足、$Q$ 为平面上其他点，则 $PQ^2=PH^2+HQ^2\ge PH^2$，故垂线段最短。长方体边长 $a,b,c>0$ 时，先求底面对角线 $\sqrt{a^2+b^2}$，再在含体对角线的矩形截面应用勾股定理，得体对角线 $\sqrt{a^2+b^2+c^2}$。两次所用直角必须来自线面关系，而非示意图外观。

\section{棱柱、棱锥与旋转体}
棱柱具有两个平行全等底面，侧棱平行等长；棱锥由一个底面多边形和一个不在其平面内的顶点组成。圆柱、圆锥可由矩形、直角三角形绕恰当边旋转形成。高是底面之间或顶点到底面所在平面的垂直距离，母线和斜高一般不是高。

直圆柱半径 $r$、高 $h$ 的侧面展开为边长 $2\pi r,h$ 的矩形，侧面积为 $2\pi rh$；直圆锥底半径 $r$、母线长 $\ell$ 的侧面展开为半径 $\ell$、弧长 $2\pi r$ 的扇形，侧面积为 $\pi r\ell$。球面是到球心距离为 $r$ 的点集，球体还包含内部。长度、面积、体积在比例 $k>0$ 的相似缩放下分别乘以 $k,k^2,k^3$。

\section{体积公式与球的度量}
采用通常体积的可加性、相似缩放规律与截面比较原理：两个立体夹在相同高度的平行平面间，若每个高度的截面积相等，则体积相等。该原理需要体积对逼近的相容性，本书作为度量基础使用，不把有限几张截面图当作证明。

底面积为 $B$、高为 $h$ 的柱体体积为 $Bh$。先注意同高同底面积的锥体在相同相对高度处截面等面积，所以截面比较保证它们等体积。把三棱柱 $ABC$--$A^{\prime}B^{\prime}C^{\prime}$ 分为四面体 $ABCC^{\prime}$、$ABB^{\prime}C^{\prime}$、$AA^{\prime}B^{\prime}C^{\prime}$：前两者共享底面 $ABC^{\prime}$，顶点 $C,B^{\prime}$ 关于该底面具有相等距离，因为侧面平行四边形的两对角线互相平分；后两者对公共底面 $AB^{\prime}C^{\prime}$ 同理等高。因此三者等体积，各为柱体的三分之一，得到锥体体积 $Bh/3$。一般底面按三角形分割后由可加性推广。平行于底面截得的锥台，两底面积为 $B_1,B_2$、高为 $h$，由相似锥体体积相减得
\[
 V=\frac h3(B_1+B_2+\sqrt{B_1B_2}).
\]

半径 $r$ 的半球在离球心平面高度 $z\in[0,r]$ 处截面积为 $\pi(r^2-z^2)$。半径、高都为 $r$ 的圆柱挖去顶点在底面中心、顶面为底的圆锥，其同高度截面积也为 $\pi r^2-\pi z^2$。因此半球体积为 $\pi r^3-\pi r^3/3=2\pi r^3/3$，整球体积为 $4\pi r^3/3$。

球面面积为 $4\pi r^2$。这里引用球面带面积定理：高度为 $h$ 的球面带面积为 $2\pi rh$；其严格证明需用内接圆台侧面积的极限定义曲面面积，不由体积公式在纯代数意义下自动推出。取整球高度 $2r$ 即得结果，相关逼近细节可在已有的圆与球度量专题中继续学习。

\section{例题、习题与学习检查}
\begin{example}
直圆锥的底半径为 $3$、高为 $4$，求母线、表面积和体积。
\end{example}
\begin{solve}
轴截面的一半是直角三角形，所以母线为 $5$。侧面积为 $15\pi$，加上底面积 $9\pi$，表面积为 $24\pi$；体积为 $\pi\cdot3^2\cdot4/3=12\pi$。面积使用平方单位，体积使用立方单位，两类结果不能混比。
\end{solve}
\begin{enumerate}
\item 长方体边长为 $2,3,6$，求体对角线。答：$7$。
\item 两相似球半径比为 $2:3$，表面积与体积比分别是多少？答：$4:9$ 与 $8:27$。
\item 高为 $H$ 的锥体在离顶点 $H/2$ 处作平行底面的截面，顶部小锥体占总体积多少？答：$1/8$。
\end{enumerate}

\part{三角、坐标与复数}

\chapter{角与三角函数}
\label{chap:em-trigonometry}
三角函数把角度与长度、坐标联系起来。单位圆定义统一了锐角、负角与超过一周的角，也使周期、奇偶性和基本恒等式成为同一几何结构的不同表现。

\section{有向角、弧度与单位圆}
以逆时针旋转为正，顺时针为负，同一终边可对应相差整周的角。\keyterm{弧度}（radian）把角度定义为所截有向弧长与半径之比，一周为 $2\pi$，故 $180^\circ=\pi$ 弧度。除特别标注角度符号外，本书三角表达式中的角均用弧度。

半径为 $r>0$、非负圆心角 $0\le\theta\le2\pi$ 的普通扇形，其弧长与面积为
\[
 s=r\theta,\qquad S=\frac12r^2\theta.
\]
这里 $s$ 是非负几何长度；对带符号的旋转角，方向与弧长大小应分别处理。多次绕行所积累的弧长也不等于终边间较短圆弧的长度。

\section{三角函数的单位圆定义}
\begin{definition}
从 $(1,0)$ 绕原点旋转角 $x$ 到单位圆上的点 $P$，其坐标写为 $(\cos x,\sin x)$，分别定义\keyterm{余弦}（cosine）和\keyterm{正弦}（sine）。当 $\cos x\ne0$ 时，定义\keyterm{正切}（tangent）$\tan x=\sin x/\cos x$。
\end{definition}
同理 $\cot x=\cos x/\sin x$、$\sec x=1/\cos x$、$\csc x=1/\sin x$，均要求分母非零。正弦与余弦在整个实轴定义；正切排除 $x=\pi/2+k\pi$，余切排除 $x=k\pi$，其中 $k\in\mathbb Z$。

对锐角，把相应直角三角形按斜边缩放到 $1$，即可看到 $\sin x$ 为对边与斜边之比、$\cos x$ 为邻边与斜边之比、$\tan x$ 为对边与邻边之比。相似性保证这些比值与三角形大小无关。

\section{定义域、符号与基本恒等式}
单位圆方程给出
\begin{equation}\label{eq:em-trig-pythagorean}
 \sin^2x+\cos^2x=1.
\end{equation}
分别除以允许非零的 $\cos^2x$ 或 $\sin^2x$，得到 $1+\tan^2x=\sec^2x$、$1+\cot^2x=\csc^2x$。记号 $\sin^2x$ 表示 $(\sin x)^2$，不是 $\sin(x^2)$。

四个象限的正弦符号依次为 $+,+,-,-$，余弦为 $+,-,-,+$，正切由二者之商确定。坐标轴上的值直接按坐标计算。等边三角形作高与等腰直角三角形分别给出
\[
 \sin\frac\pi6=\frac12,\quad\cos\frac\pi6=\frac{\sqrt3}2,\quad
 \sin\frac\pi4=\cos\frac\pi4=\frac{\sqrt2}2,\quad
 \sin\frac\pi3=\frac{\sqrt3}2,\quad\cos\frac\pi3=\frac12.
\]
若由 $\sin x$ 求 $\cos x$，平方关系只能给出 $\cos x=\pm\sqrt{1-\sin^2x}$，符号必须由角的范围确定。

\section{诱导关系、周期与图像}
单位圆关于横轴、原点及直线 $y=x$ 的对称给出
\[
 \sin(-x)=-\sin x,\quad \cos(-x)=\cos x,\quad
 \sin(x+\pi)=-\sin x,\quad\cos(x+\pi)=-\cos x,
\]
以及 $\sin(\pi/2-x)=\cos x$。转过整周后位置不变，所以正弦、余弦以 $2\pi$ 为周期，正切以 $\pi$ 为周期；它们分别没有更小的正周期，可由零点位置与相邻区间的符号验证。

在半周 $[-\pi/2,\pi/2]$ 上，单位圆纵坐标随角严格增加，故正弦严格递增；余弦在 $[0,\pi]$ 上严格递减。结合对称与周期便能确定全部单调区间。图像中正弦与余弦的值域为 $[-1,1]$，正切在每个 $(k\pi-\pi/2,k\pi+\pi/2)$ 上严格递增并覆盖实轴。一般周期函数可能没有最小正周期，例如常函数的每个正数都是周期。

\section{图像参数与反三角函数选读}
对 $y=A\sin(\omega x+\varphi)+d$，若 $A\ne0,\omega\ne0$，振幅为 $|A|$，最小正周期为 $2\pi/|\omega|$，中线为 $y=d$。相位应通过 $\omega(x+\varphi/\omega)$ 解释水平位移，不能把 $\varphi$ 本身直接当作平移长度。若 $A=0$ 或 $\omega=0$，函数退化为常量，前述最小周期公式不适用。

把正弦限制在 $[-\pi/2,\pi/2]$，余弦限制在 $[0,\pi]$，正切限制在 $(-\pi/2,\pi/2)$，它们均单射，分别定义反正弦 $\arcsin$、反余弦 $\arccos$ 和反正切 $\arctan$。反正弦、反余弦输入在 $[-1,1]$，反正切输入可为任意实数。主值只是指定区间内的一个角，不能代表三角方程的全部解。

\section{例题、习题与学习检查}
\begin{example}
已知 $x$ 的终边在第二象限，且 $\sin x=3/5$，求 $\cos x,\tan x$。
\end{example}
\begin{solve}
由式 \eqref{eq:em-trig-pythagorean} 有 $\cos^2x=16/25$，第二象限余弦为负，所以 $\cos x=-4/5$，进而 $\tan x=-3/4$。象限条件决定了平方根的符号。
\end{solve}
\begin{enumerate}
\item 将 $150^\circ$ 化为弧度并求其正弦、余弦。答：$5\pi/6,1/2,-\sqrt3/2$。
\item 求 $2\sin(3x-\pi)+1$ 的振幅、周期和值域。答：$2,2\pi/3,[-1,3]$。
\item 求 $\arcsin(\sin(5\pi/6))$。答：$\pi/6$，反函数只反转所选单射分支。
\end{enumerate}

\chapter{三角恒等变换、方程与解三角形}
\label{chap:em-triangle-solving}
三角恒等式不宜孤立记忆。和差角公式产生倍角与半角，改变加减方式又产生和积互化；这些公式与解三角形共同构成三角计算的核心。

\section{和差角公式与推导链}
\begin{theorem}[和差角公式]
对实数 $u,v$，
\begin{align}
 \cos(u\pm v)&=\cos u\cos v\mp\sin u\sin v,\label{eq:em-cos-add}\\
 \sin(u\pm v)&=\sin u\cos v\pm\cos u\sin v.\label{eq:em-sin-add}
\end{align}
\end{theorem}
\begin{proof}
单位圆上角 $u,v$ 对应点间的距离平方，一方面由坐标与勾股定理为 $2-2(\cos u\cos v+\sin u\sin v)$；另一方面旋转 $-v$ 保持距离，变为角 $u-v$ 对应点到 $(1,0)$ 的距离平方 $2-2\cos(u-v)$。相等后得到余弦差角公式，令 $v$ 换成 $-v$ 得和角式。再以 $\sin t=\cos(\pi/2-t)$ 得到正弦公式，论证没有使用后续向量数量积。
\end{proof}
在 $\cos u,\cos v,\cos(u\pm v)$ 均非零时相除，得到
\[
 \tan(u\pm v)=\frac{\tan u\pm\tan v}{1\mp\tan u\tan v}.
\]
这些分母条件是公式的组成部分，而非可忽略的附注。

\section{倍角、半角与和积互化}
令 $u=v=x$ 并使用平方关系，可得
\[
 \sin2x=2\sin x\cos x,\qquad
 \cos2x=\cos^2x-\sin^2x=1-2\sin^2x=2\cos^2x-1.
\]
于是 $\sin^2(x/2)=(1-\cos x)/2$、$\cos^2(x/2)=(1+\cos x)/2$；开方后的正负号仍由半角所在象限决定。把和角与差角式相加或相减，还得到
\[
 2\sin u\cos v=\sin(u+v)+\sin(u-v),\quad
 2\cos u\cos v=\cos(u+v)+\cos(u-v),
\]
\[
 2\sin u\sin v=\cos(u-v)-\cos(u+v).
\]
反过来令 $u=(a+b)/2,v=(a-b)/2$，便得到和差化积，例如
$\sin a+\sin b=2\sin((a+b)/2)\cos((a-b)/2)$，
$\cos a-\cos b=-2\sin((a+b)/2)\sin((a-b)/2)$。

若 $t=\tan(x/2)$ 有定义，半角式进一步给出
\[
 \sin x=\frac{2t}{1+t^2},\qquad
 \cos x=\frac{1-t^2}{1+t^2},\qquad
 \tan x=\frac{2t}{1-t^2}
\]
，最后一式还要求 $t\ne\pm1$。这种代换把三角式化为有理式，但排除了 $x=(2k+1)\pi$，解原方程时须另查这些角。

\section{三角恒等式、最值与变形策略}
恒等变形常以统一函数、降低次数或制造乘积为目标。对不同时为零的实数 $a,b$，令 $R=\sqrt{a^2+b^2}$，选择 $\varphi$ 满足 $\cos\varphi=a/R,\sin\varphi=b/R$，则
\[
 a\sin x+b\cos x=R\sin(x+\varphi).
\]
因此全实轴上的最大、最小值为 $\pm R$，且能够在相应相位取得。用 $\tan\varphi=b/a$ 单独确定角可能丢失象限，尤其在 $a=0$ 时根本无法这样除。

若 $x$ 被限制在区间，必须检查相位是否可达。例如 $\sin x+\cos x=\sqrt2\sin(x+\pi/4)$ 在 $[0,\pi/2]$ 上最大值为 $\sqrt2$，在 $x=\pi/4$ 取得；最小值为 $1$，在两个端点取得，而不是全实轴下的 $-\sqrt2$。

\section{三角方程与不等式}
对 $|c|\le1$，记 $\alpha=\arcsin c,\beta=\arccos c$，则
\[
 \sin x=c\Longleftrightarrow x=\alpha+2k\pi\ \text{或}\ x=\pi-\alpha+2k\pi,
 \quad
 \cos x=c\Longleftrightarrow x=\pm\beta+2k\pi.
\]
对任意实数 $c$，$\tan x=c$ 的解为 $x=\arctan c+k\pi$，其中 $k\in\mathbb Z$。极端值时两族表达可能重复，不产生额外解。

例如 $\sin2x=\sin x$ 化为 $\sin x(2\cos x-1)=0$，所以 $x=k\pi$ 或 $x=2k\pi\pm\pi/3$；若直接除以 $\sin x$，会漏掉整族解。三角不等式先在一个周期上确定区间，再平移：$\sin x\ge1/2$ 的解为所有 $[\pi/6+2k\pi,5\pi/6+2k\pi]$ 的并。

\section{正弦定理、余弦定理与解三角形}
设非退化三角形的边 $a,b,c$ 分别对内角 $\alpha,\beta,\gamma$，外接圆半径为 $R$。用弦长与所对圆周角关系得到\keyterm{正弦定理}（sine rule）
\[
 \frac a{\sin\alpha}=\frac b{\sin\beta}=\frac c{\sin\gamma}=2R.
\]
对锐角可把弦的一半放入以 $R$ 为斜边的直角三角形；钝角用补角正弦相同归约。把两边从共同顶点投影到适当坐标轴，再用勾股定理，可得\keyterm{余弦定理}（cosine rule）
\begin{equation}\label{eq:em-cosine-rule}
 c^2=a^2+b^2-2ab\cos\gamma,\qquad
 S=\frac12ab\sin\gamma.
\end{equation}
后一个面积式来自高为 $b\sin\gamma$；钝角时垂足在延长线上，但高仍为正。若记半周长 $s=(a+b+c)/2$，由 $16S^2=4a^2b^2-(a^2+b^2-c^2)^2$ 分解可得 Heron 公式 $S=\sqrt{s(s-a)(s-b)(s-c)}$。

三边已知先用余弦定理；两角一边已知先补第三角再用正弦定理；两边与夹角已知先求第三边。若已知 $a,b,\alpha$，则 $\sin\beta=b\sin\alpha/a$。该比值大于 $1$ 无解，等于 $1$ 时只有直角候选，小于 $1$ 时有两个互补角候选；每个候选还须满足 $\alpha+\beta<\pi$，因此可能得到零、一或两个三角形。

\section{例题、习题与学习检查}
\begin{example}
已知 $a=1,b=\sqrt3,\alpha=\pi/6$，求所有可能的三角形。
\end{example}
\begin{solve}
$\sin\beta=\sqrt3/2$，所以 $\beta=\pi/3$ 或 $2\pi/3$。两者都满足内角和限制，对应 $\gamma=\pi/2$ 或 $\pi/6$。再由 $c=\sin\gamma/\sin\alpha$ 得 $c=2$ 或 $1$，故确有两解。
\end{solve}
\begin{enumerate}
\item 求 $\sin(\pi/12)$。答：$(\sqrt6-\sqrt2)/4$。
\item 解 $\cos2x=0$。答：$x=\pi/4+k\pi/2$，$k\in\mathbb Z$。
\item 两边长为 $3,4$，夹角为 $60^\circ$，求第三边与面积。答：$\sqrt{13}$ 与 $3\sqrt3$。
\end{enumerate}

\chapter{平面向量、直线与圆的坐标方法}
\label{chap:em-vectors}
向量表达方向与位移，坐标把这些关系化为实数运算。解析几何的关键不在于写出更多方程，而在于选择坐标并证明代数条件与原几何条件等价。

\section{向量、运算与线性组合}
\keyterm{向量}（vector）具有大小和方向，可由有向线段表示；同向且等长的有向线段表示同一自由向量。零向量大小为零，不指定方向。沿折线依次移动定义向量加法，$\overrightarrow{AB}+\overrightarrow{BC}=\overrightarrow{AC}$；平行四边形法则给出等价表示。

实数乘法 $\lambda\boldsymbol u$ 把长度乘以 $|\lambda|$，正数保持方向、负数反向，零数给出零向量。形如 $\lambda\boldsymbol u+\mu\boldsymbol v$ 的式子称为线性组合。对非零 $\boldsymbol v$，$\boldsymbol u$ 与它共线当且仅当 $\boldsymbol u=\lambda\boldsymbol v$；若允许零向量，代数共线判定仍可成立，但不能据此谈零向量的几何方向。

\section{坐标、数量积与度量}
选择相互垂直的单位向量 $\boldsymbol e_1,\boldsymbol e_2$ 后，每个平面向量唯一表示为 $u_1\boldsymbol e_1+u_2\boldsymbol e_2$，记为 $(u_1,u_2)$。加法、数乘按坐标进行，勾股定理给出 $|\boldsymbol u|=\sqrt{u_1^2+u_2^2}$，两点距离为 $\sqrt{(x_2-x_1)^2+(y_2-y_1)^2}$。

对两个非零向量，\keyterm{数量积}（dot product）定义为 $\boldsymbol u\cdot\boldsymbol v=|\boldsymbol u||\boldsymbol v|\cos\theta$，其中 $\theta\in[0,\pi]$ 是夹角；含零向量时定义为零。由式 \eqref{eq:em-cosine-rule} 计算 $|\boldsymbol u-\boldsymbol v|^2$，并与坐标展开比较，得到
\[
 \boldsymbol u\cdot\boldsymbol v=u_1v_1+u_2v_2.
\]
因此数量积具有分配律。非零向量垂直当且仅当数量积为零，平行的坐标判定为 $u_1v_2-u_2v_1=0$。沿非零 $\boldsymbol v$ 的有向投影长度为 $(\boldsymbol u\cdot\boldsymbol v)/|\boldsymbol v|$，符号记录投影方向。

\section{分点、直线方程与位置关系}
在线段 $AB$ 内，若 $AP:PB=m:n$，其中 $m,n>0$，则
\[
 P=\frac{nA+mB}{m+n}
\]
按坐标逐项计算，因为 $\overrightarrow{AP}=\frac m{m+n}\overrightarrow{AB}$。中点是 $m=n$ 的情形；三角形重心坐标为三个顶点坐标的平均，代入每条中线的 $2:1$ 分点式即可核验。

过点 $P_0=(x_0,y_0)$、方向向量 $\boldsymbol v=(v_1,v_2)\ne0$ 的直线写为 $P=P_0+t\boldsymbol v$，$t\in\mathbb R$。与其垂直的非零法向量为 $(A,B)$ 时，一般式为 $Ax+By+C=0$，要求 $A^2+B^2>0$。$B\ne0$ 时才可写成斜率式 $y=kx+b$，$B=0$ 对应竖直线。两直线的法向量成比例时须再检查常数项以区分平行与重合，否则联立有唯一交点。

\section{距离、夹角与面积的坐标计算}
设点 $P=(x_0,y_0)$，直线 $\ell:Ax+By+C=0$。令
\[
 H=P-\frac{Ax_0+By_0+C}{A^2+B^2}(A,B).
\]
代入直线式得到 $H\in\ell$，且 $\overrightarrow{PH}$ 平行于法向量，故 $H$ 为垂足。于是
\begin{equation}\label{eq:em-point-line}
 d(P,\ell)=\frac{|Ax_0+By_0+C|}{\sqrt{A^2+B^2}}.
\end{equation}
两直线夹角通常取锐角或直角，方向向量 $\boldsymbol u,\boldsymbol v$ 给出 $\cos\theta=|\boldsymbol u\cdot\boldsymbol v|/(|\boldsymbol u||\boldsymbol v|)$，适用于含竖直线的情形。

若三角形从同一顶点出发的两边向量为 $\boldsymbol u,\boldsymbol v$，由 $S=|\boldsymbol u||\boldsymbol v|\sin\theta/2$ 以及二维平方和恒等式可得
\[
 S=\frac12|u_1v_2-u_2v_1|.
\]
绝对值使面积与顶点排列方向无关；表达式为零则三点共线，不能再视为非退化三角形。

\section{圆的方程与直线圆的位置}
圆心 $(a,b)$、半径 $r>0$ 的圆由距离定义得到 $(x-a)^2+(y-b)^2=r^2$。对 $x^2+y^2+Dx+Ey+F=0$ 配方，半径平方为 $(D^2+E^2)/4-F$；它为正才是非退化实圆，为零时仅有一点，为负时没有实点。

直线与圆相离、相切或相交，分别对应圆心到直线的距离大于、等于或小于半径。这也可通过联立方程的实解数验证。若 $P=(x_0,y_0)$ 在圆上，切线法向量为 $(x_0-a,y_0-b)$，故切线式为
\[
 (x_0-a)(x-a)+(y_0-b)(y-b)=r^2.
\]
点必须在圆上，才能把该式解释为所给点处的切线。

\section{例题、习题与学习检查}
\begin{example}
求直线 $3x+4y=10$ 与圆 $x^2+y^2=4$ 的位置关系和切点。
\end{example}
\begin{solve}
圆心到直线的距离为 $10/5=2$，恰等于半径，所以相切。用垂足公式或沿法向量计算，切点为 $(6/5,8/5)$；其坐标同时满足两个方程。
\end{solve}
\begin{enumerate}
\item 求过 $(1,2)$ 且法向量为 $(2,-1)$ 的直线。答：$2x-y=0$。
\item 三角形顶点为 $(0,0),(4,0),(1,3)$，求面积和重心。答：$6$ 与 $(5/3,1)$。
\item 判断 $x^2+y^2+2x-4y+6=0$ 是否有实点。答：无，配方得 $(x+1)^2+(y-2)^2=-1$。
\end{enumerate}

\chapter{圆锥曲线}
\label{chap:em-conics}
曲线方程应当是几何条件的等价表达。距离定义往往含有根式，推导方程需要平方，但反向检验必须恢复被平方隐藏的符号和位置限制。

\section{从几何定义到曲线方程}
\keyterm{圆锥曲线}（conic section）包括圆、椭圆、双曲线和抛物线。本章从焦点距离及焦点准线关系出发讨论非退化实曲线；任意二元二次方程的完整分类需要进一步的代数工具。

求轨迹的一般步骤是选择坐标、表示距离、保留定义域和位置条件、消去根式、验证所得方程上的每一点满足原条件。若原题排除了某条线或某个点，消元后仍须排除。坐标轴应尽量沿对称轴设置，使对称性体现在方程中。

\section{椭圆的定义与标准方程}
给定焦点 $F_1=(-c,0),F_2=(c,0)$，$c>0$。到两焦点距离之和为 $2a$ 的点集称为\keyterm{椭圆}（ellipse），要求 $a>c$。设两距离为 $r_1,r_2$，由
\[
 r_1+r_2=2a,\qquad r_1^2-r_2^2=4cx
\]
得到 $r_1=a+cx/a,\ r_2=a-cx/a$。代入距离平方并整理，得到
\begin{equation}\label{eq:em-ellipse}
 \frac{x^2}{a^2}+\frac{y^2}{b^2}=1,\qquad b^2=a^2-c^2>0.
\end{equation}
反之，满足该式的点有 $|x|\le a$，故 $a\pm cx/a\ge a-c>0$；它们的平方分别等于两距离平方，所以确为两距离，和为 $2a$。这样完成了必要性与充分性。

椭圆关于两坐标轴和原点对称，顶点为 $(\pm a,0),(0,\pm b)$，长、短轴长为 $2a,2b$。离心率 $e=c/a\in(0,1)$ 描述扁平程度，焦点重合的 $c=0$ 特例成为圆。以右焦点为例，$PF_2=a-cx/a=e(a^2/c-x)$，后者等于 $e$ 倍的点到准线 $x=a^2/c$ 的距离，从而也得到焦点准线表述。

\section{双曲线的定义、两支与渐近线}
同样取两焦点，距离差绝对值为 $2a$ 的点集称为\keyterm{双曲线}（hyperbola），其中 $0<a<c$。对右支先写 $r_1-r_2=2a$，由平方差得 $r_1+r_2=2cx/a$，所以 $r_1=cx/a+a,r_2=cx/a-a$，推导
\begin{equation}\label{eq:em-hyperbola}
 \frac{x^2}{a^2}-\frac{y^2}{b^2}=1,\qquad b^2=c^2-a^2.
\end{equation}
右支满足 $x\ge a$，因此候选距离均为正，反向平方检验成立；左支由关于纵轴的对称获得。方程自动给出 $|x|\ge a$，两顶点为 $(\pm a,0)$，离心率 $e=c/a>1$。

渐近线为 $y=\pm(b/a)x$。例如右上支 $y=(b/a)\sqrt{x^2-a^2}$，在 $x\ge a$ 时
\[
 \frac ba x-y=\frac{ab}{x+\sqrt{x^2-a^2}}\longrightarrow0
 \quad(x\longrightarrow+\infty).
\]
因而点到对应直线的距离也趋于零；这给出了“靠近渐近线”的准确含义，不表示曲线最终与直线重合。

\section{抛物线的定义与参数约定}
到定点与不经过它的定直线距离相等的点集称为\keyterm{抛物线}（parabola）。取焦点 $(p,0)$、准线 $x=-p$，其中 $p>0$，距离关系为
\[
 \sqrt{(x-p)^2+y^2}=|x+p|,
\]
平方整理得到
\begin{equation}\label{eq:em-parabola}
 y^2=4px.
\end{equation}
反向看，方程推出 $x\ge0$，故 $x+p>0$，距离平方相等能恢复原等式，没有多余分支。顶点为原点，对称轴为横轴，开口向右。其他方向通过交换坐标或变号处理；本书参数始终采用 $4p$ 的系数，故焦点横坐标为 $p$，不能混用其他记号体系。

抛物线的焦点准线距离比恒为 $1$，对应离心率 $1$。写成 $x=y^2/(4p)$ 后，它也是以 $y$ 为输入的二次函数图像；是否能写成 $y=f(x)$ 则取决于是否选择上下分支。

\section{直线交点、切线与综合轨迹}
将参数直线 $P=P_0+t\boldsymbol v$（$\boldsymbol v\ne0$）代入曲线，若得到 $At^2+Bt+C=0$ 且 $A\ne0$，就由判别式判断两个不同交点、重合交点或无实交点。两不同交点间弦长为 $|\boldsymbol v|\sqrt{B^2-4AC}/|A|$，中点参数为 $-B/(2A)$。若 $A=0$，必须改作一次或常数方程讨论。

在非退化圆锥曲线上，相切的代数判定是交点对应二重根，而不能只数不同交点。例如横轴 $y=0$ 与 $y^2=4px$ 只有一个交点，却是对称轴；顶点切线是 $x=0$。通过点 $(x_0,y_0)$ 的切线公式分别为
\[
 \frac{x_0x}{a^2}+\frac{y_0y}{b^2}=1,\qquad
 \frac{x_0x}{a^2}-\frac{y_0y}{b^2}=1,\qquad
 y_0y=2p(x+x_0),
\]
依次对应式 \eqref{eq:em-ellipse}、\eqref{eq:em-hyperbola}、\eqref{eq:em-parabola}，且给定点必须满足相应曲线方程。这些式子可由直线代入后要求对应参数的常数项和一次项均为零、只剩二次项来验证，无需导数。双曲线中平行于渐近线的方向可能使二次项消失，单个交点不代表二重根。

\section{例题、习题与学习检查}
\begin{example}
椭圆 $x^2/9+y^2/4=1$ 的焦点与离心率是什么？求它与横线 $y=1$ 所截弦长。
\end{example}
\begin{solve}
$a=3,b=2,c=\sqrt5$，焦点为 $(\pm\sqrt5,0)$，离心率为 $\sqrt5/3$。代入 $y=1$ 得 $x^2/9=3/4$，两交点横坐标为 $\pm3\sqrt3/2$，弦长为 $3\sqrt3$。
\end{solve}
\begin{enumerate}
\item 求 $x^2/4-y^2/5=1$ 的焦点和渐近线。答：$(\pm3,0)$，$y=\pm(\sqrt5/2)x$。
\item 求 $y^2=8x$ 的焦点、准线及点 $(2,4)$ 处切线。答：$(2,0)$、$x=-2$、$y=x+2$。
\item 两焦点距离为 $6$，距离和为 $4$，能否形成椭圆？答：不能，距离和不小于焦点间距。
\end{enumerate}

\chapter{复数及其几何意义}
\label{chap:em-complex}
方程 $x^2+1=0$ 在实数范围无解，促使我们扩展数系。复数扩展保留四则运算结构，并把乘法解释为平面上的缩放与旋转；它并不是给实数轴补上另一种大小顺序。

\section{复数、相等与四则运算}
\keyterm{复数}（complex number）写为 $z=x+iy$，其中 $x,y\in\mathbb R$，$i^2=-1$。也可严格地把它定义为实数对 $(x,y)$，规定加法为逐项相加、乘法为 $(x,y)(u,v)=(xu-yv,xv+yu)$，从而保证该表示和运算一致。实部为 $x$，虚部为实数 $y$，不是 $iy$。

复数相等当且仅当实部、虚部分别相等。展开并用 $i^2=-1$ 可得乘法公式；对 $u+iv\ne0$，
\[
 \frac{x+iy}{u+iv}=
 \frac{(xu+yv)+i(yu-xv)}{u^2+v^2}.
\]
分母为正，因此每个非零复数都有唯一倒数。不能给复数建立与实数有序域规则相容的顺序，否则平方都应非负，而 $i^2=-1<0$ 造成矛盾。

\section{共轭、模与复平面}
定义\keyterm{共轭}（complex conjugate）$\overline z=x-iy$ 和\keyterm{模}（modulus）$|z|=\sqrt{x^2+y^2}$。有
\[
 z\overline z=|z|^2,\qquad
 \overline{z+w}=\overline z+\overline w,\qquad
 \overline{zw}=\overline z\,\overline w.
\]
将乘积与其共轭相乘得 $|zw|^2=|z|^2|w|^2$，取非负平方根得到 $|zw|=|z||w|$。

在复平面中 $z$ 对应 $(x,y)$，加法是位移相加，共轭是关于实轴反射，$|z-w|$ 是两点距离。平面三角不等式给出 $|z+w|\le|z|+|w|$；含零值时取等，非零时取等要求二者同方向。$|z-a|=r$（$r>0$）表示以 $a$ 为圆心的圆，$|z-a|=|z-b|$（$a\ne b$）表示两点连线的垂直平分线。

\section{三角形式、辐角与乘除法}
若 $z\ne0$，记 $r=|z|>0$，则
\[
 z=r(\cos\theta+i\sin\theta).
\]
角 $\theta$ 称为\keyterm{辐角}（argument），所有可能值相差 $2k\pi$；需要单值时采用主辐角区间 $(-\pi,\pi]$。零复数没有确定方向，因此不定义辐角。

由和差角公式，两个非零复数相乘时模相乘、辐角相加；相除时模相除、辐角相减。乘以 $i$ 是逆时针旋转 $\pi/2$，乘以正数是等比例缩放，乘以负数再附加半周旋转。主辐角相加后可能落在规定区间外，须加减整周重新归入区间。

\section{复数乘方、开方与单位根}
对非零复数及整数 $n$，反复乘法和取倒数给出 de Moivre 公式
\[
 [r(\cos\theta+i\sin\theta)]^n
 =r^n(\cos n\theta+i\sin n\theta).
\]
设 $w=\rho(\cos\varphi+i\sin\varphi)\ne0$，正整数 $n$ 次方程 $z^n=w$ 的全部解为
\[
 z_k=\rho^{1/n}\left(\cos\frac{\varphi+2k\pi}{n}
       +i\sin\frac{\varphi+2k\pi}{n}\right),
 \qquad k=0,\ldots,n-1.
\]
必要性来自模的 $n$ 次方与辐角同余，充分性由代入验证；这 $n$ 个角模 $2\pi$ 互不相同，所以解也互异。$w=0$ 时唯一解为 $z=0$。

$z^n=1$ 的根称为 $n$ 次单位根，均匀分布在单位圆上。若 $\zeta\ne1$ 且 $\zeta^n=1$，则等比和给出 $1+\zeta+\cdots+\zeta^{n-1}=0$；这也反映了规则旋转配置的向量平衡。

\section{二次方程与几何应用}
实系数二次方程在 $\Delta<0$ 时，两根为
\[
 z=\frac{-b\pm i\sqrt{-\Delta}}{2a},\qquad a\ne0,
\]
互为共轭。更一般地，实系数多项式满足 $f(\overline z)=\overline{f(z)}$，故非实根成共轭对；这只说明已有根的对称性，不是一般多项式根存在性的证明。

围绕复点 $a$ 旋转 $\theta$ 的变换可写为
\[
 z'=a+(\cos\theta+i\sin\theta)(z-a).
\]
先减去中心、旋转位移、再加回中心，说明运算顺序的几何意义。复数坐标使旋转简洁，但长度关系仍需通过模表达，不能把复数本身当作有序长度。

\section{例题、习题与学习检查}
\begin{example}
求 $z^3=8$ 的全部复数解。
\end{example}
\begin{solve}
模为 $2$，辐角可取 $0,2\pi/3,4\pi/3$，故三根为 $2,-1+i\sqrt3,-1-i\sqrt3$。后三角表示中的整周重复不再产生新根；实数求根只保留了第一根。
\end{solve}
\begin{enumerate}
\item 化简 $(1+i)/(1-i)$。答：$i$。
\item 求 $z^2-2z+5=0$ 的根。答：$1\pm2i$。
\item 将 $z=2+i$ 绕 $a=1$ 逆时针旋转 $90^\circ$。答：$z'=1+i(1+i)=i$。
\end{enumerate}

\part{随机、数据与综合应用}

\chapter{有限概率模型}
\label{chap:em-probability}
概率为随机现象指定可计算的模型。随机并不意味着每种描述都等可能，概率计算必须先说明基本结果是什么，以及为什么采用所给权重。本章只讨论有限样本空间，避免引入尚未建立的测度理论。

\section{随机试验、样本空间与事件}
\keyterm{样本空间}（sample space）$\Omega$ 是一次试验所有可能基本结果组成的有限集合，\keyterm{事件}（event）是其子集。两次掷硬币可用有序对 $(\text{正},\text{反})$ 等表示；“至少一次正面”包含三个有序结果，而“正面次数为零、一、二”这三类大小不同，不能因为有三类就各赋概率 $1/3$。

并、交、补分别表达至少一个发生、同时发生和不发生。$A\cap B=\varnothing$ 称为互斥，$B=A^c$ 称为对立；对立还要求并为整个样本空间，因此比互斥更强。基本结果的区分应与试验记录和概率赋值相匹配。

\section{概率、古典概型与加法公式}
给每个 $\omega\in\Omega$ 指定 $p(\omega)\ge0$ 且 $\sum_\omega p(\omega)=1$，定义
\[
 \mathbb P(A)=\sum_{\omega\in A}p(\omega).
\]
于是 $\mathbb P(\varnothing)=0,\mathbb P(\Omega)=1$，互斥事件概率可相加。有限等可能模型才满足 $\mathbb P(A)=|A|/|\Omega|$。对一般事件，分割为不重叠部分得到
\[
 \mathbb P(A\cup B)=\mathbb P(A)+\mathbb P(B)-\mathbb P(A\cap B),
 \qquad \mathbb P(A^c)=1-\mathbb P(A).
\]
例如从 $5$ 个不同球中等可能无放回抽取 $2$ 个，若其中 $2$ 个为红球，则两红概率为 $\binom22/\binom52=1/10$。采用无序结果是合法的，因为每个无序对对应相同数量的有序抽取。

\section{条件概率、独立与互斥}
在 $\mathbb P(B)>0$ 时，定义\keyterm{条件概率}（conditional probability）
\[
 \mathbb P(A\mid B)=\frac{\mathbb P(A\cap B)}{\mathbb P(B)}.
\]
它把注意范围缩到 $B$，再重新归一化。由定义得乘法公式 $\mathbb P(A\cap B)=\mathbb P(B)\mathbb P(A\mid B)$，并可逐步扩展到多次条件选择。

事件 $A,B$ \keyterm{独立}（independent）是指 $\mathbb P(A\cap B)=\mathbb P(A)\mathbb P(B)$；当 $\mathbb P(B)>0$ 时等价于条件信息不改变 $A$ 的概率。两个正概率互斥事件不独立，因为交集概率为零而概率乘积为正。多个事件联合独立要求每个子族的交概率都等于相应乘积，仅两两独立不充分。例如两个独立公平比特为 $X,Y$，事件 $X=0,Y=0,X=Y$ 两两独立，但三者同时发生概率为 $1/4$，不等于 $(1/2)^3$。

\section{全概率与贝叶斯公式}
设 $B_1,\ldots,B_m$ 两两互斥且并为 $\Omega$，并有 $\mathbb P(B_i)>0$。按来源分割 $A$ 得
\begin{equation}\label{eq:em-total-probability}
 \mathbb P(A)=\sum_{i=1}^m\mathbb P(B_i)\mathbb P(A\mid B_i).
\end{equation}
若 $\mathbb P(A)>0$，再用条件概率定义得到\keyterm{贝叶斯公式}（Bayes' formula）
\[
 \mathbb P(B_j\mid A)=
 \frac{\mathbb P(B_j)\mathbb P(A\mid B_j)}
 {\sum_i\mathbb P(B_i)\mathbb P(A\mid B_i)}.
\]
这是从观察结果更新来源概率的规则，而非交换条件符号。分母必须包含全部可能来源，条件概率也不是来源的先验概率。

\section{有限随机变量、期望与方差}
\keyterm{随机变量}（random variable）是函数 $X:\Omega\to\mathbb R$，$X$ 表示随机对象，$x$ 表示某个取值。其分布列为 $p_X(x)=\mathbb P(X=x)$，只需列出实际可能值。定义
\[
 \mathbb E[X]=\sum_x x\,p_X(x),\qquad
 \operatorname{Var}(X)=\mathbb E[(X-\mathbb E[X])^2]
 =\mathbb E[X^2]-\mathbb E[X]^2.
\]
所有和有限，故无可积性问题。按基本结果逐项相加可证明 $\mathbb E[aX+bY]=a\mathbb E[X]+b\mathbb E[Y]$，这不要求独立；独立时才可进一步保证乘积期望分解以及方差相加。

取值 $0,1$ 且 $\mathbb P(X=1)=p$ 的变量称为 Bernoulli 变量，它的期望为 $p$、方差为 $p(1-p)$。$n$ 次联合独立、成功概率同为 $p$ 的试验成功总数 $S$ 服从\keyterm{二项分布}（binomial distribution）。当 $0<p<1$ 时，
\[
 \mathbb P(S=k)=\binom nk p^k(1-p)^{n-k},\quad 0\le k\le n,
 \qquad \mathbb E[S]=np,\quad\operatorname{Var}(S)=np(1-p).
\]
概率式先指定哪 $k$ 次成功再求和；期望由线性性，方差由独立性推出。$p=0,1$ 时分别为恒等于 $0,n$ 的退化情形，无需给 $0^0$ 赋值。期望相同的变量可能具有不同方差，因此期望不完整描述结果波动。

\section{例题、习题与学习检查}
\begin{example}
等可能选择两个盒子之一，第一盒含两红一白，第二盒含一红三白。从选中盒中等可能取一球，观察为红球，求来自第一盒的概率。
\end{example}
\begin{solve}
两来源先验概率均为 $1/2$，红球条件概率分别为 $2/3,1/4$，故红球总概率为 $1/3+1/8=11/24$。第一盒且红球的概率为 $1/3$，后验概率为 $(1/3)/(11/24)=8/11$。
\end{solve}
\begin{enumerate}
\item 三次独立公平掷币恰有两次正面的概率是多少？答：$3/8$。
\item 公平六面骰点数的期望是多少？答：$7/2$；期望未必是可能取值。
\item 若 $Y=X$，其中 $X$ 是公平 Bernoulli 变量，比较 $\operatorname{Var}(X+Y)$ 与 $\operatorname{Var}(X)+\operatorname{Var}(Y)$。答：分别为 $1$ 与 $1/2$，缺少独立性时不能直接相加。
\end{enumerate}

\chapter{统计描述与数据解释}
\label{chap:em-statistics}
统计描述把有限数据压缩为可解释的信息，统计推断则进一步使用抽样模型认识总体。压缩必然损失细节，因此应根据问题选择指标，同时检查数据来源、观测单位和抽样方式。本章数值例子均为教学自拟数据。

\section{总体、样本与数据来源}
\keyterm{总体}（population）是研究涉及的全部对象，\keyterm{样本}（sample）是实际观测到的部分对象；每个对象上的特征称为变量。抽样前 $X_1,\ldots,X_n$ 是随机变量，观测后得到具体值 $x_1,\ldots,x_n$。不含未知总体参数、仅由样本及已知量计算的量称为\keyterm{统计量}（statistic）。

普查试图覆盖总体，抽样则只观察部分单位。简单随机抽样、分层抽样等需有明确选择规则；便利样本和自愿回答样本可能系统性偏向某类对象。样本量增大能减少某些随机波动，却不能自动消除选择偏差、漏报或测量偏差。描述“参加调查的人”的结果，不应未经论证扩展为“所有人”的结论。

\section{数据类型、频数与图表}
类别数据表示种类，数值数据表示可进行有意义算术运算的量；给类别编码为 $1,2,3$ 并不自动赋予平均值实际含义。某类出现次数为频数，除以总观测数得到频率，各互斥完备类别的频率和为 $1$。

条形图比较类别，直方图表示分组数值数据。若第 $j$ 组宽度为 $w_j>0$、频率为 $f_j$，应以 $f_j/w_j$ 为高度，使面积代表频率；组宽不同时直接比较柱高会误导。散点图保留两个数值变量的配对关系。分组后一般不能精确恢复原数据，利用组中值计算均值是一种近似，不是原样本均值的恒等变形。

\section{集中趋势与平均数的选择}
样本均值为 $\overline x=n^{-1}\sum_i x_i$。带非负权重 $w_i$ 且总权重为正时，加权均值为 $\sum_iw_ix_i/\sum_iw_i$。把若干组汇总时，组均值应按组样本量加权，而不是无条件取组均值的简单平均。

将数据排序为 $x_{(1)}\le\cdots\le x_{(n)}$，中位数在奇数样本量时取中间一项，在偶数时取中间两项平均；众数为出现次数最多的值，可以有多个。均值利用全部数值，易受极端值影响；中位数主要反映位置顺序。例如自拟数据 $1,2,2,3,12$ 的均值为 $4$，中位数与众数均为 $2$，差异揭示右侧大值的影响，而不说明某一种指标普遍优越。

均值还有最小二乘意义：对任意实数 $a$，
\[
 \sum_i(x_i-a)^2=\sum_i(x_i-\overline x)^2+n(a-\overline x)^2,
\]
交叉项因 $\sum_i(x_i-\overline x)=0$ 消失。因此均值是使平方偏差和最小的唯一常数。绝对偏差和则由中位数最小化：把两端数据成对，$|a-u|+|a-v|$ 在 $a\in[u,v]$ 时最小，所有这些区间的交给出中位位置。

\section{离散程度、标准化与样本方差}
极差为最大值减最小值。对 $n\ge2$，本书的四分位数约定是：先取总体中位数，奇数样本量时去掉该中间项，再分别取上下半数据的中位数作为 $Q_1,Q_3$；四分位距为 $Q_3-Q_1$。其他软件可能采用插值约定，小样本比较时应核对定义。

描述当前数据离散程度可用
\[
 v_n=\frac1n\sum_i(x_i-\overline x)^2,\qquad
 \sqrt{v_n},
\]
分别为以 $n$ 作分母的方差和标准差。$n\ge2$ 时校正样本方差为 $s^2=(n-1)^{-1}\sum_i(x_i-\overline x)^2$。它的统计意义依赖抽样：若 $X_i$ 独立同分布且取值有限、总体均值为 $\mu$、方差为 $\sigma^2$，则
\[
 \sum_i(X_i-\overline X)^2=\sum_i(X_i-\mu)^2-n(\overline X-\mu)^2.
\]
取期望并用 $\operatorname{Var}(\overline X)=\sigma^2/n$，得到期望为 $(n-1)\sigma^2$，故 $\mathbb E[s^2]=\sigma^2$；以 $n$ 为分母则期望为 $(n-1)\sigma^2/n$。此处 $s^2$ 指观测前的随机统计量，其数值在观测后才确定。标准化 $(x_i-\overline x)/s$ 还要求 $s>0$。

\section{相关、线性拟合与推断边界}
对配对数据 $(x_i,y_i)$，记中心化值 $u_i=x_i-\overline x,v_i=y_i-\overline y$。当 $\sum_i u_i^2$ 与 $\sum_i v_i^2$ 均非零时，样本相关系数为
\[
 r=\frac{\sum_i u_iv_i}{\sqrt{(\sum_i u_i^2)(\sum_i v_i^2)}}.
\]
由柯西不等式，$|r|\le1$；等号意味着中心化数据成比例，即全部点落在某条非水平、非竖直直线上。$r=0$ 只表示这种线性指标为零，不保证没有非线性关系；相关也不能单独推出因果。

拟合直线 $y=\alpha+\beta x$ 的最小二乘目标是使 $\sum_i(y_i-\alpha-\beta x_i)^2$ 最小。令 $D=\sum_i u_i^2>0$、$N=\sum_i u_iv_i$，配方得到
\[
 \sum_i(y_i-\alpha-\beta x_i)^2
 =\sum_i v_i^2-\frac{N^2}{D}
  +D\left(\beta-\frac ND\right)^2
  +n(\alpha+\beta\overline x-\overline y)^2.
\]
因此唯一最优解为 $\widehat\beta=N/D,\widehat\alpha=\overline y-\widehat\beta\overline x$。这是一个确定性的拟合结论，不需要先假设误差正态；若要作置信区间或因果解释，则还需额外模型。所有 $x_i$ 相同的 $D=0$ 情形不能唯一确定斜率。

\section{例题、习题与学习检查}
\begin{example}
对自拟数据 $(0,1),(1,2),(2,2)$ 求最小二乘直线。
\end{example}
\begin{solve}
$\overline x=1,\overline y=5/3$，中心化后的 $D=2,N=1$，所以 $\widehat\beta=1/2,\widehat\alpha=7/6$。拟合直线为 $y=7/6+x/2$，三个残差依次为 $-1/6,1/3,-1/6$，其和为零、平方和为 $1/6$。拟合并不保证每个观测点落在直线上。
\end{solve}
\begin{enumerate}
\item 数据 $2,3,3,4,8$ 的均值、中位数与 $v_n$ 是多少？答：$4,3,22/5$。
\item 一组 $10$ 人平均值为 $60$，另一组 $30$ 人平均值为 $80$，合并平均值是多少？答：$75$。
\item 若把所有数据加常数 $c$，均值和方差怎样变化？答：均值加 $c$，方差不变；乘常数 $c$ 时方差乘 $c^2$。
\end{enumerate}

\chapter{数学建模、综合方法与后续学习}
\label{chap:em-modeling}
综合应用要求在不同数学表示之间作选择，并把计算结果送回原问题检验。一个模型可以在数学上完全正确，却因假设不合实际而失去解释力；因此变量、约束和适用范围与求解过程同样重要。

\section{从问题到变量、约束与目标}
\keyterm{数学建模}（mathematical modeling）包括识别对象、提出假设、定义变量、建立关系、求解与检验。未知量是需要确定的值，参数描述外部条件，约束规定允许的取值，目标则说明怎样比较方案。方程描述相等关系，不等式描述容量、范围或预算，整数约束描述不可分割的数量。

例如行程关系 $s=vt$ 依赖速度在所讨论时段恒定；混合物守恒需要所计物质在过程中不被产生或损失。不同单位应先统一，出现负长度、概率超过 $1$ 或应为整数的数量出现小数时，应检查模型与筛选，而非机械接受代数输出。模型假设改变时，原解法可能仍正确，但已经回答另一个问题。

\section{函数、方程与不等式的综合使用}
用长度为 $L>0$ 的围栏围成靠直墙的矩形，墙边无需围栏。令垂直墙边长为 $x$、平行墙边长为 $y$，则 $2x+y=L$、$0<x<L/2$，面积为
\[
 S=x(L-2x)=\frac{L^2}{8}-2\left(x-\frac L4\right)^2.
\]
所以最大面积为 $L^2/8$，在 $x=L/4,y=L/2$ 取得。也可把约束看成 $2x+y=L$，对两正数 $2x,y$ 用均值不等式得到 $2xy\le L^2/4$。两方法分别从顶点与平衡取等解释同一结构。

若 $L=19$ 且两边长须为正整数，连续最优点 $x=19/4$ 不可行。配方说明只需比较允许整数中最接近它的 $4,5$：对应 $(x,y)=(4,11),(5,9)$，面积为 $44,45$，故整数最优方案是 $(5,9)$。不能独立地把连续解各坐标四舍五入，因为那可能破坏 $2x+y=19$。

\section{几何、三角与坐标的综合使用}
考虑水平地面上的竖直目标，忽略观测者高度。在同一直线上靠近与远离目标的两站间距为 $d>0$，仰角分别为 $\alpha,\beta$，满足 $0<\beta<\alpha<\pi/2$。设近站到目标底部距离为 $x>0$、目标高为 $h$，则
\[
 h=x\tan\alpha=(x+d)\tan\beta.
\]
消元得
\[
 x=\frac{d\tan\beta}{\tan\alpha-\tan\beta},\qquad
 h=\frac{d\tan\alpha\tan\beta}{\tan\alpha-\tan\beta}.
\]
若 $d=10,\alpha=45^\circ,\beta=30^\circ$，则 $h=5(\sqrt3+1)$。分母为正与近站仰角较大相一致；若两角非常接近，小测量误差可能显著影响结果，应保留角度误差范围。观测者高度、地面坡度或目标倾斜存在时，须相应修改几何关系。

\section{随机模拟、数据与模型检验}
给定有限概率模型、试验次数 $N$ 和目标事件 $A$，模拟按模型独立生成结果，以指示变量 $I_j$ 记录第 $j$ 次是否发生 $A$，输出经验频率 $\widehat p=N^{-1}\sum_j I_j$。设模型概率为 $p$，则
\[
 \mathbb E[\widehat p]=p,\qquad
 \operatorname{Var}(\widehat p)=\frac{p(1-p)}N.
\]
独立性用于方差计算，伪随机程序也只有在被合理用作该模型的近似生成器时，才支持这种解释。

对任意 $\varepsilon>0$，由于偏差超过 $\varepsilon$ 时平方偏差至少为 $\varepsilon^2$，
\[
 \varepsilon^2\mathbb P(|\widehat p-p|\ge\varepsilon)
 \le\mathbb E[(\widehat p-p)^2]
 =\frac{p(1-p)}N,
\]
故超出误差阈值的概率至多为 $1/(4N\varepsilon^2)$。这是一个可能较松的理论保证，不是某次模拟必然达到的精度。例如 $N=10000,\varepsilon=0.05$ 时上界为 $0.01$。模拟可以发现错误和形成猜测，却不能替代普遍成立的数学证明。

\section{通向高等数学的概念桥梁}
初等数学已经产生了高等数学的主要问题：有限和如何走向无穷和，平均变化如何走向瞬时变化，直线逼近如何刻画曲线，方程组的消元怎样推广到高维。数学分析研究极限、连续、微分与积分，线性代数研究向量与线性映射，概率论把有限模型推广到更一般的随机空间。

以自然对数的底数为例，令 $e_n=(1+1/n)^n$。对 $n$ 个 $1+1/n$ 和一个 $1$ 使用算术--几何平均不等式，可得
\[
 e_n<\left(1+\frac1{n+1}\right)^{n+1}=e_{n+1}.
\]
又二项式定理给出 $e_n\le\sum_{k=0}^n1/k!$。当 $k\ge2$ 时 $k!\ge2\cdot3^{k-2}$，所以对 $n\ge2$，
\[
 e_n\le2+\frac12\sum_{j=0}^{n-2}3^{-j}<\frac{11}4.
\]
故数列递增有界，其极限定义为 $e$，并有 $2<e<3$，从而底为 $e$ 的对数 $\ln x$ 有合法定义。证明 $(e^x)'=e^x$ 等性质还需要建立导数及相应极限定理，本书没有因定义了 $e$ 就提前假定这些结论。这个例子说明，数列、计数、不等式与完备性如何共同支持一个新的数学对象。

\section{综合习题、表达训练与结课自测}
\begin{example}
一个模型给出面积上界 $M$，但所有达到等号的参数都违反原约束。能否宣布最大面积为 $M$？
\end{example}
\begin{solve}
不能。已证的只是每个合法方案的面积不超过 $M$；还需证明某个合法方案达到它，或重新寻找可达的更小上界。若能构造合法方案使面积任意接近 $M$，则 $M$ 是上确界，但仍可能没有最大值。连续模型与整数模型、开区间与闭区间之间的差异常由此产生。
\end{solve}
\begin{enumerate}
\item 靠墙矩形围栏长为 $20$，求最大面积和边长。答：$50$，垂直墙边长 $5$、平行墙边长 $10$。
\item 在上述独立模拟模型中，使误差超过 $0.1$ 的概率上界不大于 $0.05$，次数至少应满足什么条件？答：$N\ge500$；这是充分保证，不一定是最小实际所需次数。
\item 用同一道几何问题比较综合法、三角法与坐标法，逐项说明它们使用的对象、假设与取等或存在条件。评价解答应看论证与适用范围，而非只看计算步数。
\end{enumerate}

\appendix

\chapter{符号、公式与方法索引}
\label{app:em-notation}
本附录用于定位工具。公式不能脱离条件单独使用，具体证明和例题应返回对应章节阅读。

\section{符号与中英文术语}
$\mathbb N$ 包含零，$\mathbb N_{>0}$ 表示正整数；$\mathbb Z,\mathbb Q,\mathbb R,\mathbb C$ 依次为整数、有理数、实数和复数集。$A\subseteq B$ 表示集合包含，$x\in A$ 表示元素关系；$|A|$ 在有限集合语境中是基数（cardinality），$|x|$ 对实数是绝对值，对复数是模。$\gcd$ 表示最大公因数，$\operatorname{lcm}$ 表示最小公倍数。

$f^{-1}$ 表示单射函数在值域上的反函数，不表示倒数。$a_n$ 是数列第 $n$ 项，$S_n$ 是前 $n$ 项和，$\binom nk$ 表示组合数。$\boldsymbol u\cdot\boldsymbol v$ 为数量积，$\overline z$ 为复共轭；统计中的 $\overline x$ 则表示样本均值，含义由所定义对象确定。$\mathbb P,\mathbb E,\operatorname{Var}$ 分别表示概率、期望和方差，随机变量用大写字母，观测值通常用小写。

\section{常用公式及其条件}
代数求解可查二次方程公式 \eqref{eq:em-quadratic}、因式定理 \eqref{eq:em-factor} 和 Bézout 等式 \eqref{eq:em-bezout}；求和可查等比和 \eqref{eq:em-geometric-sum} 与二项式定理 \eqref{eq:em-binomial}。不等式的变量范围和取等条件见第\ref{chap:em-inequality-methods}章，特别是式 \eqref{eq:em-cauchy}。

三角推导的起点为平方关系 \eqref{eq:em-trig-pythagorean} 与和差角公式 \eqref{eq:em-cos-add}、\eqref{eq:em-sin-add}，解三角形用式 \eqref{eq:em-cosine-rule}。坐标距离用式 \eqref{eq:em-point-line}，三类圆锥曲线的参数约定分别见式 \eqref{eq:em-ellipse}、\eqref{eq:em-hyperbola}、\eqref{eq:em-parabola}。概率分来源计算用式 \eqref{eq:em-total-probability}，使用前须检查分割完备性。

\section{问题结构与方法检索}
多项式零点先看因式分解和余式，分式符号先找全部零点与禁用点，含绝对值表达式先用距离或分段。定和定积约束优先检验均值不等式和配方，重复累积先识别等差、等比或可平移的递推。几何中的等长与比例联系全等、相似与点幂，位置和方向可转化为向量或坐标。

有限选择先明确顺序、重复及对象等价关系，再用加法、乘法或容斥；存在性可尝试抽屉原理，不可达性可寻找不变量。概率先建样本空间和权重，统计先查抽样口径与单位。每个方法都应以满足其适用条件为入口，关键词相似本身不构成使用依据。

\chapter{分层自测与专题阅读}
\label{app:em-study}
学习检查应同时覆盖概念、计算、证明和解释。能够识别一个熟悉公式，并不等于能够说明它的条件；能够跟随一个证明，也不等于能够独立补足关键步骤。

\section{分阶段自测与复习路径}
第一阶段检查集合包含与元素关系、充分必要条件、绝对值、整除及定义域。第二阶段检查方程变形是否等价、函数是否单射、参数何时退化、数列公式是否覆盖初值。第三阶段检查不等式等号能否达到、计数是否重复、有限证明中所选对象是否存在。

几何阶段检查图形中的关系是否来自条件、全等相似的对应是否一致、轨迹是否双向验证、作图是否遗漏退化情况；三角与坐标阶段检查角度单位、分支、法向量和曲线参数约定。概率统计阶段检查模型是否等可能、独立性是否给出、样本结论能否扩展到总体。遇到概念性错误应回到定义，遇到运算性错误应逐步验算，遇到证明性错误应寻找缺失依据或反例。

\section{综合自测及答案要点}
\begin{enumerate}
\item 解 $\dfrac{x^2-1}{x-1}\ge0$。答：$[-1,1)\cup(1,+\infty)$，先保留 $x\ne1$ 再化简。
\item 正数 $x,y$ 满足 $x+y=10$，求 $x^2+y^2$ 的最小值。答：$50$，由 $(x-y)^2\ge0$，在 $x=y=5$ 取得。
\item 求从 $7$ 个不同对象中选出 $3$ 个的方法数，并求同时选中两个指定对象的概率，假设所有三元子集等可能。答：总数 $35$，有利数 $5$，概率 $1/7$。
\item 三角形三边为 $5,5,6$，求面积、内切圆和外接圆半径。答：高为 $4$，面积 $12$；由 $S=rs$ 得内切半径 $3/2$，由 $S=abc/(4R)$ 得外接半径 $25/8$。前一关系来自内心分成的三个三角形，后一关系结合正弦定理与面积式。
\item 解复数方程 $z^2=-4$。答：$\pm2i$；实数范围内则无解，先确定数系。
\item 给出一个样本均值正确但总体解释可能错误的情形。要点：例如仅调查自愿参加者却把结果扩展到所有对象，错误在抽样代表性而非均值公式。
\end{enumerate}

\section{现有专题与进一步阅读}
现有均值、初等不等式与 Jensen 不等式专题适合在第\ref{chap:em-inequality-methods}章之后阅读，其中涉及积分或一般随机变量的推广需补足相应分析与概率基础。裂项公式和双阶乘专题分别衔接第\ref{chap:em-sequences}、\ref{chap:em-counting}章；三角公式专题衔接第\ref{chap:em-triangle-solving}章；圆与球的度量专题可用于进一步追踪长度、面积和体积逼近的严格依据。

继续学习时，数学分析承接实数完备性、数列与函数，线性代数承接方程组和向量，初等数论与抽象代数深化整除及代数结构，概率论与数理统计扩展有限模型和抽样推断。已有知识的价值在于明确哪些工具已经证明、哪些条件不可省略，以及新问题究竟超出了现有理论的哪一步。

\end{document}
